% 高一04_位育中学_周末作业1.tex —— 高一上数学「2026-2027 年位育中学高一上周末作业 1」（试卷版/答案版）
% 试卷版：xelatex 高一04_位育中学_周末作业1.tex
% 答案版：xelatex "\def\isanswer{1}% 高一04_位育中学_周末作业1.tex —— 高一上数学「2026-2027 年位育中学高一上周末作业 1」（试卷版/答案版）
% 试卷版：xelatex 高一04_位育中学_周末作业1.tex
% 答案版：xelatex "\def\isanswer{1}% 高一04_位育中学_周末作业1.tex —— 高一上数学「2026-2027 年位育中学高一上周末作业 1」（试卷版/答案版）
% 试卷版：xelatex 高一04_位育中学_周末作业1.tex
% 答案版：xelatex "\def\isanswer{1}% 高一04_位育中学_周末作业1.tex —— 高一上数学「2026-2027 年位育中学高一上周末作业 1」（试卷版/答案版）
% 试卷版：xelatex 高一04_位育中学_周末作业1.tex
% 答案版：xelatex "\def\isanswer{1}\input{高一04_位育中学_周末作业1.tex}"
% 紧凑版：xelatex "\def\iscompact{1}\input{高一04_位育中学_周末作业1.tex}"
%
% 原始材料是微信公众号图文（「魔都数学加油站」，作者署名「破晓」），正文为 3 张 PNG 页
% （1080×1527）；卷面上的粉色斜水印「魔都数学加油站」属水印，未录入。原文本身就是一份
% **讲评版**——题下直接印出【答案】或【解析】，本稿把这些答案与解析移到答案版，试卷版即
% 一张干净的空卷。
%
% 与原卷的排版差异（均已在 README「说明」中交代）：
%   ① 原文自**第 3 题**起（第 1、2 题未在该文发布），源码里以 \setcounter{enumi}{2} 起编；
%   ② 原卷本就印有分节标题「一、填空题」「二、选择题」「三、解答题」（已在原图上逐页放大
%      核对过），本稿照原卷分节，只把标题改排成全稿统一的「大节标题」样式；
%   ③ 原卷的实数集、整数集符号**正体与粗体混用**——第 3 题作正体 $N$、$Z$，第 6 题作粗体
%      $\mathbf{R}$（均逐处放大核对过），本稿按全稿惯例统一写作 $\N$、$\Z$、$\R$；
%   ④ 原卷第 6、7 题的 $\subset$ 不带下划线（即真包含，与带下划线的 $\subseteq$ 在卷面上
%      逐处放大比对后确认不同），本稿按原符号保留，并按真包含口径解答（正文内有 `%` 注）.

\documentclass[a4paper,11pt]{article}
\usepackage{common}

\begin{document}

\examtitlest{2026--2027 学年度位育中学高一上周末作业 1}{集合与常用逻辑用语}

\bigsec{一、填空题}

\begin{enumerate}
\setcounter{enumi}{2}

% 注：本行照原卷录入——原卷作「$y=\dfrac{12}{3-x}\in N,\,x\in Z$」，即「$y=\frac{12}{3-x}$ 是自然数、$x$ 是整数」；
%     本稿曾误写成「$x\in\mathbb{N},\,x\in\mathbb{Z}$」（把 $\in\mathbb{N}$ 挂到了 $x$ 上），
%     那样只能得到 $\{4,6,12\}$，与本卷给的答案 $\{1,2,3,4,6,12\}$ 冲突。
%     2026-09-15 回原图核对时已按原卷订正。
\item 用列举法表示集合 $A=\left\{y\,\middle|\,y=\dfrac{12}{3-x}\in\mathbb{N},\,x\in\mathbb{Z}\right\}=$ \blankline[2em].

\ans{$\{1,2,3,4,6,12\}$}

\item 已知集合 $A=\{2a,\,2a^2+a\}$，若 $1\in A$，则 $a=$ \blankline[2em].

\ans{$-1$.}

\ifanswer
\solbegin
\step{因为 $1\in A$，所以 $2a=1$ 或 $2a^2+a=1$.}
\step{当 $2a=1$ 时，$a=\tfrac{1}{2}$，此时 $2a^2+a=1$，不满足集合中元素的互异性，舍去；}
\step{当 $2a^2+a=1$ 时，解得 $a=-1$ 或 $a=\tfrac{1}{2}$（舍去），此时 $A=\{-2,1\}$.}
\step{综上，$a=-1$.}
\solend

\fi
\item 关于 $x,y$ 的方程组 $\begin{cases} x+y=2 \\ x^2+y^2=2 \end{cases}$ 的解集为 \blankline[2em].

\ans{$\{(1,1)\}$}

\item 如果集合 $A=\{x\mid y=x^2,\,x\in\R\}$，$B=\{y\mid y=x^2,\,x\in\R\}$，则集合 $A$ 与 $B$ 的关系是 \blankline[2em].

% 注：本行答案照原卷作「$B\subset A$」（原卷的 $\subset$ 不带下划线，语义即真包含，与 $\subsetneq$ 等价）；
%     本稿原写作 $\subsetneq$，2026-09-15 回原图核对时按原卷记号改回 $\subset$。
\ans{$B\subset A$}

% 注：本题答案照原卷作「$A\subset B$」（原卷此处的 $\subset$ 不带下划线，即真包含）——
%     $A=\left\{\frac{m}{6}\,\middle|\,m\text{ 为奇数}\right\}$、$B=\left\{\frac{m}{6}\,\middle|\,m\in\Z\right\}$，确是 $A\subsetneq B$。
%     本稿曾误抄成「$A\supsetneq B$（或 $A=B$ 之等价条件）」——方向与结论都反了，
%     2026-09-15 回原图核对时已按原卷订正。
\item 已知 $A=\left\{x\,\middle|\,x=\dfrac{k}{3}-\dfrac{1}{2},\,k\in\Z\right\}$，$B=\left\{x\,\middle|\,x=\dfrac{k}{6}+\dfrac{2}{3},\,k\in\Z\right\}$，则集合 $A$ 与 $B$ 的关系是 \blankline[2em].

\ans{$A\subset B$}

\item 若非空集合 $\{x\mid m\leqslant x\leqslant 2m-\sqrt{m}\}$ 不存在非空真子集，则 $m$ 的取值集合为 \blankline[2em].

\ans{$\{0,1\}$}

\item 已知 $P=\{x\mid 2<x<a,\,x\in\mathbb{N}\}$，且集合 $P$ 中恰有 3 个元素，则实数 $a$ 的取值范围为 \blankline[3em]（用区间表示）.

\ans{$(5,6]$}

\item 设 $x_1,x_2,x_3,x_4$ 是 4 个正整数，从中任取 3 个数求和所得的集合为 $\{25,26,27\}$，则这 4 个数中最小的数为 \blankline[2em].

\ans{$8$}

\ifanswer
\solbegin
\step{从 4 个正整数中任取 3 个数求和后可得到 4 个和，则 4 个和值之和为 $3(x_1+x_2+x_3+x_4)$，必为 3 的倍数.}
\step{又 $27\div 3=9,\,25\div 3=8\dots1,\,26\div 3=8\dots2$，所以这 4 个和为 $25,26,27,27$.}
\step{则 $x_1+x_2+x_3+x_4=\tfrac{1}{3}(25+26+27+27)=35$.}
\step{所以 $35-25=10,\,35-26=9,\,35-27=8$，即这 4 个数分别为 $8,8,9,10$，故这 4 个数中最小的数为 8.}
\solend

\fi
\end{enumerate}

\bigsec{二、选择题}

\begin{enumerate}
\setcounter{enumi}{10}

\item 下列说法中正确的是\mbox{（\quad）.}

\keepwithprev
A. $0$ 与 $\{0\}$ 表示同一个集合；

B. 集合 $\{1,2,3\}$ 与 $\{3,2,1\}$ 是两个相同的集合；

C. 方程 $(x-1)^2(x-2)=0$ 的解集为 $\{1,1,2\}$；

D. 集合 $\{x\mid 4<x<5\}$ 可以用列举法表示.

\ans{B}

\item 下列说法正确的是\mbox{（\quad）.}

\keepwithprev
A. $\varnothing\in\{0\}$\qquad B. $0\subseteq\{0\}$\qquad C. $\varnothing\subseteq\{0,1,2\}$\qquad D. $(1,2)\supset\{1,2\}$

\ans{C}

\item 已知非空集合 $A=\{x\mid x^2-2x+m=0\}$，则集合 $A$ 中所有元素之和为\mbox{（\quad）.}

\keepwithprev
A. $1$\qquad B. $2$\qquad C. $1$ 或 $2$\qquad D. 非以上答案

\ans{C}

\item 如果集合 $S$ 具有性质：① 非空且它的元素都是正整数；② 如果 $x\in S$，那么 $\tfrac{24}{x}\in S$. 请问这样的集合 $S$ 共有\mbox{（\quad）}个.

\keepwithprev
A. $15$ 个\qquad B. $16$ 个\qquad C. $31$ 个\qquad D. $32$ 个

\ans{A}

\end{enumerate}

\bigsec{三、解答题}

\begin{enumerate}
\setcounter{enumi}{14}

\item 设集合 $A=\{x\mid x^2-ax+4=0\}$，$B=\{x\mid x^2-3x+2=0\}$，且 $A\subseteq B$，求实数 $a$ 的取值范围.

\ans{$\{a\mid -4<a\leqslant 4\}$（注：原答案如此）}

\ansspace[6]
\ansneed{7cm}
\item 已知集合 $A=\{x\mid -1<x<6\}$，$B=\{x\mid -m<x<2m+1\}$，其中 $m$ 为实常数.
\begin{enumerate}[label=(\arabic*)]
\item 若 $B\subseteq A$，求实数 $m$ 的取值范围；
% 注：本小问原卷作「求证：$A\subset B$」（不带下划线，即真包含），本稿原写作 $\subseteq$，
%     2026-09-15 回原图核对时按原卷记号订正。
\item 若 $m>3$，求证：$A\subset B$.
\end{enumerate}

\ans{（1）$m\leqslant 1$；（2）略}

\ansspace*[7]
\end{enumerate}

\end{document}
"
% 紧凑版：xelatex "\def\iscompact{1}% 高一04_位育中学_周末作业1.tex —— 高一上数学「2026-2027 年位育中学高一上周末作业 1」（试卷版/答案版）
% 试卷版：xelatex 高一04_位育中学_周末作业1.tex
% 答案版：xelatex "\def\isanswer{1}\input{高一04_位育中学_周末作业1.tex}"
% 紧凑版：xelatex "\def\iscompact{1}\input{高一04_位育中学_周末作业1.tex}"
%
% 原始材料是微信公众号图文（「魔都数学加油站」，作者署名「破晓」），正文为 3 张 PNG 页
% （1080×1527）；卷面上的粉色斜水印「魔都数学加油站」属水印，未录入。原文本身就是一份
% **讲评版**——题下直接印出【答案】或【解析】，本稿把这些答案与解析移到答案版，试卷版即
% 一张干净的空卷。
%
% 与原卷的排版差异（均已在 README「说明」中交代）：
%   ① 原文自**第 3 题**起（第 1、2 题未在该文发布），源码里以 \setcounter{enumi}{2} 起编；
%   ② 原卷本就印有分节标题「一、填空题」「二、选择题」「三、解答题」（已在原图上逐页放大
%      核对过），本稿照原卷分节，只把标题改排成全稿统一的「大节标题」样式；
%   ③ 原卷的实数集、整数集符号**正体与粗体混用**——第 3 题作正体 $N$、$Z$，第 6 题作粗体
%      $\mathbf{R}$（均逐处放大核对过），本稿按全稿惯例统一写作 $\N$、$\Z$、$\R$；
%   ④ 原卷第 6、7 题的 $\subset$ 不带下划线（即真包含，与带下划线的 $\subseteq$ 在卷面上
%      逐处放大比对后确认不同），本稿按原符号保留，并按真包含口径解答（正文内有 `%` 注）.

\documentclass[a4paper,11pt]{article}
\usepackage{common}

\begin{document}

\examtitlest{2026--2027 学年度位育中学高一上周末作业 1}{集合与常用逻辑用语}

\bigsec{一、填空题}

\begin{enumerate}
\setcounter{enumi}{2}

% 注：本行照原卷录入——原卷作「$y=\dfrac{12}{3-x}\in N,\,x\in Z$」，即「$y=\frac{12}{3-x}$ 是自然数、$x$ 是整数」；
%     本稿曾误写成「$x\in\mathbb{N},\,x\in\mathbb{Z}$」（把 $\in\mathbb{N}$ 挂到了 $x$ 上），
%     那样只能得到 $\{4,6,12\}$，与本卷给的答案 $\{1,2,3,4,6,12\}$ 冲突。
%     2026-09-15 回原图核对时已按原卷订正。
\item 用列举法表示集合 $A=\left\{y\,\middle|\,y=\dfrac{12}{3-x}\in\mathbb{N},\,x\in\mathbb{Z}\right\}=$ \blankline[2em].

\ans{$\{1,2,3,4,6,12\}$}

\item 已知集合 $A=\{2a,\,2a^2+a\}$，若 $1\in A$，则 $a=$ \blankline[2em].

\ans{$-1$.}

\ifanswer
\solbegin
\step{因为 $1\in A$，所以 $2a=1$ 或 $2a^2+a=1$.}
\step{当 $2a=1$ 时，$a=\tfrac{1}{2}$，此时 $2a^2+a=1$，不满足集合中元素的互异性，舍去；}
\step{当 $2a^2+a=1$ 时，解得 $a=-1$ 或 $a=\tfrac{1}{2}$（舍去），此时 $A=\{-2,1\}$.}
\step{综上，$a=-1$.}
\solend

\fi
\item 关于 $x,y$ 的方程组 $\begin{cases} x+y=2 \\ x^2+y^2=2 \end{cases}$ 的解集为 \blankline[2em].

\ans{$\{(1,1)\}$}

\item 如果集合 $A=\{x\mid y=x^2,\,x\in\R\}$，$B=\{y\mid y=x^2,\,x\in\R\}$，则集合 $A$ 与 $B$ 的关系是 \blankline[2em].

% 注：本行答案照原卷作「$B\subset A$」（原卷的 $\subset$ 不带下划线，语义即真包含，与 $\subsetneq$ 等价）；
%     本稿原写作 $\subsetneq$，2026-09-15 回原图核对时按原卷记号改回 $\subset$。
\ans{$B\subset A$}

% 注：本题答案照原卷作「$A\subset B$」（原卷此处的 $\subset$ 不带下划线，即真包含）——
%     $A=\left\{\frac{m}{6}\,\middle|\,m\text{ 为奇数}\right\}$、$B=\left\{\frac{m}{6}\,\middle|\,m\in\Z\right\}$，确是 $A\subsetneq B$。
%     本稿曾误抄成「$A\supsetneq B$（或 $A=B$ 之等价条件）」——方向与结论都反了，
%     2026-09-15 回原图核对时已按原卷订正。
\item 已知 $A=\left\{x\,\middle|\,x=\dfrac{k}{3}-\dfrac{1}{2},\,k\in\Z\right\}$，$B=\left\{x\,\middle|\,x=\dfrac{k}{6}+\dfrac{2}{3},\,k\in\Z\right\}$，则集合 $A$ 与 $B$ 的关系是 \blankline[2em].

\ans{$A\subset B$}

\item 若非空集合 $\{x\mid m\leqslant x\leqslant 2m-\sqrt{m}\}$ 不存在非空真子集，则 $m$ 的取值集合为 \blankline[2em].

\ans{$\{0,1\}$}

\item 已知 $P=\{x\mid 2<x<a,\,x\in\mathbb{N}\}$，且集合 $P$ 中恰有 3 个元素，则实数 $a$ 的取值范围为 \blankline[3em]（用区间表示）.

\ans{$(5,6]$}

\item 设 $x_1,x_2,x_3,x_4$ 是 4 个正整数，从中任取 3 个数求和所得的集合为 $\{25,26,27\}$，则这 4 个数中最小的数为 \blankline[2em].

\ans{$8$}

\ifanswer
\solbegin
\step{从 4 个正整数中任取 3 个数求和后可得到 4 个和，则 4 个和值之和为 $3(x_1+x_2+x_3+x_4)$，必为 3 的倍数.}
\step{又 $27\div 3=9,\,25\div 3=8\dots1,\,26\div 3=8\dots2$，所以这 4 个和为 $25,26,27,27$.}
\step{则 $x_1+x_2+x_3+x_4=\tfrac{1}{3}(25+26+27+27)=35$.}
\step{所以 $35-25=10,\,35-26=9,\,35-27=8$，即这 4 个数分别为 $8,8,9,10$，故这 4 个数中最小的数为 8.}
\solend

\fi
\end{enumerate}

\bigsec{二、选择题}

\begin{enumerate}
\setcounter{enumi}{10}

\item 下列说法中正确的是\mbox{（\quad）.}

\keepwithprev
A. $0$ 与 $\{0\}$ 表示同一个集合；

B. 集合 $\{1,2,3\}$ 与 $\{3,2,1\}$ 是两个相同的集合；

C. 方程 $(x-1)^2(x-2)=0$ 的解集为 $\{1,1,2\}$；

D. 集合 $\{x\mid 4<x<5\}$ 可以用列举法表示.

\ans{B}

\item 下列说法正确的是\mbox{（\quad）.}

\keepwithprev
A. $\varnothing\in\{0\}$\qquad B. $0\subseteq\{0\}$\qquad C. $\varnothing\subseteq\{0,1,2\}$\qquad D. $(1,2)\supset\{1,2\}$

\ans{C}

\item 已知非空集合 $A=\{x\mid x^2-2x+m=0\}$，则集合 $A$ 中所有元素之和为\mbox{（\quad）.}

\keepwithprev
A. $1$\qquad B. $2$\qquad C. $1$ 或 $2$\qquad D. 非以上答案

\ans{C}

\item 如果集合 $S$ 具有性质：① 非空且它的元素都是正整数；② 如果 $x\in S$，那么 $\tfrac{24}{x}\in S$. 请问这样的集合 $S$ 共有\mbox{（\quad）}个.

\keepwithprev
A. $15$ 个\qquad B. $16$ 个\qquad C. $31$ 个\qquad D. $32$ 个

\ans{A}

\end{enumerate}

\bigsec{三、解答题}

\begin{enumerate}
\setcounter{enumi}{14}

\item 设集合 $A=\{x\mid x^2-ax+4=0\}$，$B=\{x\mid x^2-3x+2=0\}$，且 $A\subseteq B$，求实数 $a$ 的取值范围.

\ans{$\{a\mid -4<a\leqslant 4\}$（注：原答案如此）}

\ansspace[6]
\ansneed{7cm}
\item 已知集合 $A=\{x\mid -1<x<6\}$，$B=\{x\mid -m<x<2m+1\}$，其中 $m$ 为实常数.
\begin{enumerate}[label=(\arabic*)]
\item 若 $B\subseteq A$，求实数 $m$ 的取值范围；
% 注：本小问原卷作「求证：$A\subset B$」（不带下划线，即真包含），本稿原写作 $\subseteq$，
%     2026-09-15 回原图核对时按原卷记号订正。
\item 若 $m>3$，求证：$A\subset B$.
\end{enumerate}

\ans{（1）$m\leqslant 1$；（2）略}

\ansspace*[7]
\end{enumerate}

\end{document}
"
%
% 原始材料是微信公众号图文（「魔都数学加油站」，作者署名「破晓」），正文为 3 张 PNG 页
% （1080×1527）；卷面上的粉色斜水印「魔都数学加油站」属水印，未录入。原文本身就是一份
% **讲评版**——题下直接印出【答案】或【解析】，本稿把这些答案与解析移到答案版，试卷版即
% 一张干净的空卷。
%
% 与原卷的排版差异（均已在 README「说明」中交代）：
%   ① 原文自**第 3 题**起（第 1、2 题未在该文发布），源码里以 \setcounter{enumi}{2} 起编；
%   ② 原卷本就印有分节标题「一、填空题」「二、选择题」「三、解答题」（已在原图上逐页放大
%      核对过），本稿照原卷分节，只把标题改排成全稿统一的「大节标题」样式；
%   ③ 原卷的实数集、整数集符号**正体与粗体混用**——第 3 题作正体 $N$、$Z$，第 6 题作粗体
%      $\mathbf{R}$（均逐处放大核对过），本稿按全稿惯例统一写作 $\N$、$\Z$、$\R$；
%   ④ 原卷第 6、7 题的 $\subset$ 不带下划线（即真包含，与带下划线的 $\subseteq$ 在卷面上
%      逐处放大比对后确认不同），本稿按原符号保留，并按真包含口径解答（正文内有 `%` 注）.

\documentclass[a4paper,11pt]{article}
\usepackage{common}

\begin{document}

\examtitlest{2026--2027 学年度位育中学高一上周末作业 1}{集合与常用逻辑用语}

\bigsec{一、填空题}

\begin{enumerate}
\setcounter{enumi}{2}

% 注：本行照原卷录入——原卷作「$y=\dfrac{12}{3-x}\in N,\,x\in Z$」，即「$y=\frac{12}{3-x}$ 是自然数、$x$ 是整数」；
%     本稿曾误写成「$x\in\mathbb{N},\,x\in\mathbb{Z}$」（把 $\in\mathbb{N}$ 挂到了 $x$ 上），
%     那样只能得到 $\{4,6,12\}$，与本卷给的答案 $\{1,2,3,4,6,12\}$ 冲突。
%     2026-09-15 回原图核对时已按原卷订正。
\item 用列举法表示集合 $A=\left\{y\,\middle|\,y=\dfrac{12}{3-x}\in\mathbb{N},\,x\in\mathbb{Z}\right\}=$ \blankline[2em].

\ans{$\{1,2,3,4,6,12\}$}

\item 已知集合 $A=\{2a,\,2a^2+a\}$，若 $1\in A$，则 $a=$ \blankline[2em].

\ans{$-1$.}

\ifanswer
\solbegin
\step{因为 $1\in A$，所以 $2a=1$ 或 $2a^2+a=1$.}
\step{当 $2a=1$ 时，$a=\tfrac{1}{2}$，此时 $2a^2+a=1$，不满足集合中元素的互异性，舍去；}
\step{当 $2a^2+a=1$ 时，解得 $a=-1$ 或 $a=\tfrac{1}{2}$（舍去），此时 $A=\{-2,1\}$.}
\step{综上，$a=-1$.}
\solend

\fi
\item 关于 $x,y$ 的方程组 $\begin{cases} x+y=2 \\ x^2+y^2=2 \end{cases}$ 的解集为 \blankline[2em].

\ans{$\{(1,1)\}$}

\item 如果集合 $A=\{x\mid y=x^2,\,x\in\R\}$，$B=\{y\mid y=x^2,\,x\in\R\}$，则集合 $A$ 与 $B$ 的关系是 \blankline[2em].

% 注：本行答案照原卷作「$B\subset A$」（原卷的 $\subset$ 不带下划线，语义即真包含，与 $\subsetneq$ 等价）；
%     本稿原写作 $\subsetneq$，2026-09-15 回原图核对时按原卷记号改回 $\subset$。
\ans{$B\subset A$}

% 注：本题答案照原卷作「$A\subset B$」（原卷此处的 $\subset$ 不带下划线，即真包含）——
%     $A=\left\{\frac{m}{6}\,\middle|\,m\text{ 为奇数}\right\}$、$B=\left\{\frac{m}{6}\,\middle|\,m\in\Z\right\}$，确是 $A\subsetneq B$。
%     本稿曾误抄成「$A\supsetneq B$（或 $A=B$ 之等价条件）」——方向与结论都反了，
%     2026-09-15 回原图核对时已按原卷订正。
\item 已知 $A=\left\{x\,\middle|\,x=\dfrac{k}{3}-\dfrac{1}{2},\,k\in\Z\right\}$，$B=\left\{x\,\middle|\,x=\dfrac{k}{6}+\dfrac{2}{3},\,k\in\Z\right\}$，则集合 $A$ 与 $B$ 的关系是 \blankline[2em].

\ans{$A\subset B$}

\item 若非空集合 $\{x\mid m\leqslant x\leqslant 2m-\sqrt{m}\}$ 不存在非空真子集，则 $m$ 的取值集合为 \blankline[2em].

\ans{$\{0,1\}$}

\item 已知 $P=\{x\mid 2<x<a,\,x\in\mathbb{N}\}$，且集合 $P$ 中恰有 3 个元素，则实数 $a$ 的取值范围为 \blankline[3em]（用区间表示）.

\ans{$(5,6]$}

\item 设 $x_1,x_2,x_3,x_4$ 是 4 个正整数，从中任取 3 个数求和所得的集合为 $\{25,26,27\}$，则这 4 个数中最小的数为 \blankline[2em].

\ans{$8$}

\ifanswer
\solbegin
\step{从 4 个正整数中任取 3 个数求和后可得到 4 个和，则 4 个和值之和为 $3(x_1+x_2+x_3+x_4)$，必为 3 的倍数.}
\step{又 $27\div 3=9,\,25\div 3=8\dots1,\,26\div 3=8\dots2$，所以这 4 个和为 $25,26,27,27$.}
\step{则 $x_1+x_2+x_3+x_4=\tfrac{1}{3}(25+26+27+27)=35$.}
\step{所以 $35-25=10,\,35-26=9,\,35-27=8$，即这 4 个数分别为 $8,8,9,10$，故这 4 个数中最小的数为 8.}
\solend

\fi
\end{enumerate}

\bigsec{二、选择题}

\begin{enumerate}
\setcounter{enumi}{10}

\item 下列说法中正确的是\mbox{（\quad）.}

\keepwithprev
A. $0$ 与 $\{0\}$ 表示同一个集合；

B. 集合 $\{1,2,3\}$ 与 $\{3,2,1\}$ 是两个相同的集合；

C. 方程 $(x-1)^2(x-2)=0$ 的解集为 $\{1,1,2\}$；

D. 集合 $\{x\mid 4<x<5\}$ 可以用列举法表示.

\ans{B}

\item 下列说法正确的是\mbox{（\quad）.}

\keepwithprev
A. $\varnothing\in\{0\}$\qquad B. $0\subseteq\{0\}$\qquad C. $\varnothing\subseteq\{0,1,2\}$\qquad D. $(1,2)\supset\{1,2\}$

\ans{C}

\item 已知非空集合 $A=\{x\mid x^2-2x+m=0\}$，则集合 $A$ 中所有元素之和为\mbox{（\quad）.}

\keepwithprev
A. $1$\qquad B. $2$\qquad C. $1$ 或 $2$\qquad D. 非以上答案

\ans{C}

\item 如果集合 $S$ 具有性质：① 非空且它的元素都是正整数；② 如果 $x\in S$，那么 $\tfrac{24}{x}\in S$. 请问这样的集合 $S$ 共有\mbox{（\quad）}个.

\keepwithprev
A. $15$ 个\qquad B. $16$ 个\qquad C. $31$ 个\qquad D. $32$ 个

\ans{A}

\end{enumerate}

\bigsec{三、解答题}

\begin{enumerate}
\setcounter{enumi}{14}

\item 设集合 $A=\{x\mid x^2-ax+4=0\}$，$B=\{x\mid x^2-3x+2=0\}$，且 $A\subseteq B$，求实数 $a$ 的取值范围.

\ans{$\{a\mid -4<a\leqslant 4\}$（注：原答案如此）}

\ansspace[6]
\ansneed{7cm}
\item 已知集合 $A=\{x\mid -1<x<6\}$，$B=\{x\mid -m<x<2m+1\}$，其中 $m$ 为实常数.
\begin{enumerate}[label=(\arabic*)]
\item 若 $B\subseteq A$，求实数 $m$ 的取值范围；
% 注：本小问原卷作「求证：$A\subset B$」（不带下划线，即真包含），本稿原写作 $\subseteq$，
%     2026-09-15 回原图核对时按原卷记号订正。
\item 若 $m>3$，求证：$A\subset B$.
\end{enumerate}

\ans{（1）$m\leqslant 1$；（2）略}

\ansspace*[7]
\end{enumerate}

\end{document}
"
% 紧凑版：xelatex "\def\iscompact{1}% 高一04_位育中学_周末作业1.tex —— 高一上数学「2026-2027 年位育中学高一上周末作业 1」（试卷版/答案版）
% 试卷版：xelatex 高一04_位育中学_周末作业1.tex
% 答案版：xelatex "\def\isanswer{1}% 高一04_位育中学_周末作业1.tex —— 高一上数学「2026-2027 年位育中学高一上周末作业 1」（试卷版/答案版）
% 试卷版：xelatex 高一04_位育中学_周末作业1.tex
% 答案版：xelatex "\def\isanswer{1}\input{高一04_位育中学_周末作业1.tex}"
% 紧凑版：xelatex "\def\iscompact{1}\input{高一04_位育中学_周末作业1.tex}"
%
% 原始材料是微信公众号图文（「魔都数学加油站」，作者署名「破晓」），正文为 3 张 PNG 页
% （1080×1527）；卷面上的粉色斜水印「魔都数学加油站」属水印，未录入。原文本身就是一份
% **讲评版**——题下直接印出【答案】或【解析】，本稿把这些答案与解析移到答案版，试卷版即
% 一张干净的空卷。
%
% 与原卷的排版差异（均已在 README「说明」中交代）：
%   ① 原文自**第 3 题**起（第 1、2 题未在该文发布），源码里以 \setcounter{enumi}{2} 起编；
%   ② 原卷本就印有分节标题「一、填空题」「二、选择题」「三、解答题」（已在原图上逐页放大
%      核对过），本稿照原卷分节，只把标题改排成全稿统一的「大节标题」样式；
%   ③ 原卷的实数集、整数集符号**正体与粗体混用**——第 3 题作正体 $N$、$Z$，第 6 题作粗体
%      $\mathbf{R}$（均逐处放大核对过），本稿按全稿惯例统一写作 $\N$、$\Z$、$\R$；
%   ④ 原卷第 6、7 题的 $\subset$ 不带下划线（即真包含，与带下划线的 $\subseteq$ 在卷面上
%      逐处放大比对后确认不同），本稿按原符号保留，并按真包含口径解答（正文内有 `%` 注）.

\documentclass[a4paper,11pt]{article}
\usepackage{common}

\begin{document}

\examtitlest{2026--2027 学年度位育中学高一上周末作业 1}{集合与常用逻辑用语}

\bigsec{一、填空题}

\begin{enumerate}
\setcounter{enumi}{2}

% 注：本行照原卷录入——原卷作「$y=\dfrac{12}{3-x}\in N,\,x\in Z$」，即「$y=\frac{12}{3-x}$ 是自然数、$x$ 是整数」；
%     本稿曾误写成「$x\in\mathbb{N},\,x\in\mathbb{Z}$」（把 $\in\mathbb{N}$ 挂到了 $x$ 上），
%     那样只能得到 $\{4,6,12\}$，与本卷给的答案 $\{1,2,3,4,6,12\}$ 冲突。
%     2026-09-15 回原图核对时已按原卷订正。
\item 用列举法表示集合 $A=\left\{y\,\middle|\,y=\dfrac{12}{3-x}\in\mathbb{N},\,x\in\mathbb{Z}\right\}=$ \blankline[2em].

\ans{$\{1,2,3,4,6,12\}$}

\item 已知集合 $A=\{2a,\,2a^2+a\}$，若 $1\in A$，则 $a=$ \blankline[2em].

\ans{$-1$.}

\ifanswer
\solbegin
\step{因为 $1\in A$，所以 $2a=1$ 或 $2a^2+a=1$.}
\step{当 $2a=1$ 时，$a=\tfrac{1}{2}$，此时 $2a^2+a=1$，不满足集合中元素的互异性，舍去；}
\step{当 $2a^2+a=1$ 时，解得 $a=-1$ 或 $a=\tfrac{1}{2}$（舍去），此时 $A=\{-2,1\}$.}
\step{综上，$a=-1$.}
\solend

\fi
\item 关于 $x,y$ 的方程组 $\begin{cases} x+y=2 \\ x^2+y^2=2 \end{cases}$ 的解集为 \blankline[2em].

\ans{$\{(1,1)\}$}

\item 如果集合 $A=\{x\mid y=x^2,\,x\in\R\}$，$B=\{y\mid y=x^2,\,x\in\R\}$，则集合 $A$ 与 $B$ 的关系是 \blankline[2em].

% 注：本行答案照原卷作「$B\subset A$」（原卷的 $\subset$ 不带下划线，语义即真包含，与 $\subsetneq$ 等价）；
%     本稿原写作 $\subsetneq$，2026-09-15 回原图核对时按原卷记号改回 $\subset$。
\ans{$B\subset A$}

% 注：本题答案照原卷作「$A\subset B$」（原卷此处的 $\subset$ 不带下划线，即真包含）——
%     $A=\left\{\frac{m}{6}\,\middle|\,m\text{ 为奇数}\right\}$、$B=\left\{\frac{m}{6}\,\middle|\,m\in\Z\right\}$，确是 $A\subsetneq B$。
%     本稿曾误抄成「$A\supsetneq B$（或 $A=B$ 之等价条件）」——方向与结论都反了，
%     2026-09-15 回原图核对时已按原卷订正。
\item 已知 $A=\left\{x\,\middle|\,x=\dfrac{k}{3}-\dfrac{1}{2},\,k\in\Z\right\}$，$B=\left\{x\,\middle|\,x=\dfrac{k}{6}+\dfrac{2}{3},\,k\in\Z\right\}$，则集合 $A$ 与 $B$ 的关系是 \blankline[2em].

\ans{$A\subset B$}

\item 若非空集合 $\{x\mid m\leqslant x\leqslant 2m-\sqrt{m}\}$ 不存在非空真子集，则 $m$ 的取值集合为 \blankline[2em].

\ans{$\{0,1\}$}

\item 已知 $P=\{x\mid 2<x<a,\,x\in\mathbb{N}\}$，且集合 $P$ 中恰有 3 个元素，则实数 $a$ 的取值范围为 \blankline[3em]（用区间表示）.

\ans{$(5,6]$}

\item 设 $x_1,x_2,x_3,x_4$ 是 4 个正整数，从中任取 3 个数求和所得的集合为 $\{25,26,27\}$，则这 4 个数中最小的数为 \blankline[2em].

\ans{$8$}

\ifanswer
\solbegin
\step{从 4 个正整数中任取 3 个数求和后可得到 4 个和，则 4 个和值之和为 $3(x_1+x_2+x_3+x_4)$，必为 3 的倍数.}
\step{又 $27\div 3=9,\,25\div 3=8\dots1,\,26\div 3=8\dots2$，所以这 4 个和为 $25,26,27,27$.}
\step{则 $x_1+x_2+x_3+x_4=\tfrac{1}{3}(25+26+27+27)=35$.}
\step{所以 $35-25=10,\,35-26=9,\,35-27=8$，即这 4 个数分别为 $8,8,9,10$，故这 4 个数中最小的数为 8.}
\solend

\fi
\end{enumerate}

\bigsec{二、选择题}

\begin{enumerate}
\setcounter{enumi}{10}

\item 下列说法中正确的是\mbox{（\quad）.}

\keepwithprev
A. $0$ 与 $\{0\}$ 表示同一个集合；

B. 集合 $\{1,2,3\}$ 与 $\{3,2,1\}$ 是两个相同的集合；

C. 方程 $(x-1)^2(x-2)=0$ 的解集为 $\{1,1,2\}$；

D. 集合 $\{x\mid 4<x<5\}$ 可以用列举法表示.

\ans{B}

\item 下列说法正确的是\mbox{（\quad）.}

\keepwithprev
A. $\varnothing\in\{0\}$\qquad B. $0\subseteq\{0\}$\qquad C. $\varnothing\subseteq\{0,1,2\}$\qquad D. $(1,2)\supset\{1,2\}$

\ans{C}

\item 已知非空集合 $A=\{x\mid x^2-2x+m=0\}$，则集合 $A$ 中所有元素之和为\mbox{（\quad）.}

\keepwithprev
A. $1$\qquad B. $2$\qquad C. $1$ 或 $2$\qquad D. 非以上答案

\ans{C}

\item 如果集合 $S$ 具有性质：① 非空且它的元素都是正整数；② 如果 $x\in S$，那么 $\tfrac{24}{x}\in S$. 请问这样的集合 $S$ 共有\mbox{（\quad）}个.

\keepwithprev
A. $15$ 个\qquad B. $16$ 个\qquad C. $31$ 个\qquad D. $32$ 个

\ans{A}

\end{enumerate}

\bigsec{三、解答题}

\begin{enumerate}
\setcounter{enumi}{14}

\item 设集合 $A=\{x\mid x^2-ax+4=0\}$，$B=\{x\mid x^2-3x+2=0\}$，且 $A\subseteq B$，求实数 $a$ 的取值范围.

\ans{$\{a\mid -4<a\leqslant 4\}$（注：原答案如此）}

\ansspace[6]
\ansneed{7cm}
\item 已知集合 $A=\{x\mid -1<x<6\}$，$B=\{x\mid -m<x<2m+1\}$，其中 $m$ 为实常数.
\begin{enumerate}[label=(\arabic*)]
\item 若 $B\subseteq A$，求实数 $m$ 的取值范围；
% 注：本小问原卷作「求证：$A\subset B$」（不带下划线，即真包含），本稿原写作 $\subseteq$，
%     2026-09-15 回原图核对时按原卷记号订正。
\item 若 $m>3$，求证：$A\subset B$.
\end{enumerate}

\ans{（1）$m\leqslant 1$；（2）略}

\ansspace*[7]
\end{enumerate}

\end{document}
"
% 紧凑版：xelatex "\def\iscompact{1}% 高一04_位育中学_周末作业1.tex —— 高一上数学「2026-2027 年位育中学高一上周末作业 1」（试卷版/答案版）
% 试卷版：xelatex 高一04_位育中学_周末作业1.tex
% 答案版：xelatex "\def\isanswer{1}\input{高一04_位育中学_周末作业1.tex}"
% 紧凑版：xelatex "\def\iscompact{1}\input{高一04_位育中学_周末作业1.tex}"
%
% 原始材料是微信公众号图文（「魔都数学加油站」，作者署名「破晓」），正文为 3 张 PNG 页
% （1080×1527）；卷面上的粉色斜水印「魔都数学加油站」属水印，未录入。原文本身就是一份
% **讲评版**——题下直接印出【答案】或【解析】，本稿把这些答案与解析移到答案版，试卷版即
% 一张干净的空卷。
%
% 与原卷的排版差异（均已在 README「说明」中交代）：
%   ① 原文自**第 3 题**起（第 1、2 题未在该文发布），源码里以 \setcounter{enumi}{2} 起编；
%   ② 原卷本就印有分节标题「一、填空题」「二、选择题」「三、解答题」（已在原图上逐页放大
%      核对过），本稿照原卷分节，只把标题改排成全稿统一的「大节标题」样式；
%   ③ 原卷的实数集、整数集符号**正体与粗体混用**——第 3 题作正体 $N$、$Z$，第 6 题作粗体
%      $\mathbf{R}$（均逐处放大核对过），本稿按全稿惯例统一写作 $\N$、$\Z$、$\R$；
%   ④ 原卷第 6、7 题的 $\subset$ 不带下划线（即真包含，与带下划线的 $\subseteq$ 在卷面上
%      逐处放大比对后确认不同），本稿按原符号保留，并按真包含口径解答（正文内有 `%` 注）.

\documentclass[a4paper,11pt]{article}
\usepackage{common}

\begin{document}

\examtitlest{2026--2027 学年度位育中学高一上周末作业 1}{集合与常用逻辑用语}

\bigsec{一、填空题}

\begin{enumerate}
\setcounter{enumi}{2}

% 注：本行照原卷录入——原卷作「$y=\dfrac{12}{3-x}\in N,\,x\in Z$」，即「$y=\frac{12}{3-x}$ 是自然数、$x$ 是整数」；
%     本稿曾误写成「$x\in\mathbb{N},\,x\in\mathbb{Z}$」（把 $\in\mathbb{N}$ 挂到了 $x$ 上），
%     那样只能得到 $\{4,6,12\}$，与本卷给的答案 $\{1,2,3,4,6,12\}$ 冲突。
%     2026-09-15 回原图核对时已按原卷订正。
\item 用列举法表示集合 $A=\left\{y\,\middle|\,y=\dfrac{12}{3-x}\in\mathbb{N},\,x\in\mathbb{Z}\right\}=$ \blankline[2em].

\ans{$\{1,2,3,4,6,12\}$}

\item 已知集合 $A=\{2a,\,2a^2+a\}$，若 $1\in A$，则 $a=$ \blankline[2em].

\ans{$-1$.}

\ifanswer
\solbegin
\step{因为 $1\in A$，所以 $2a=1$ 或 $2a^2+a=1$.}
\step{当 $2a=1$ 时，$a=\tfrac{1}{2}$，此时 $2a^2+a=1$，不满足集合中元素的互异性，舍去；}
\step{当 $2a^2+a=1$ 时，解得 $a=-1$ 或 $a=\tfrac{1}{2}$（舍去），此时 $A=\{-2,1\}$.}
\step{综上，$a=-1$.}
\solend

\fi
\item 关于 $x,y$ 的方程组 $\begin{cases} x+y=2 \\ x^2+y^2=2 \end{cases}$ 的解集为 \blankline[2em].

\ans{$\{(1,1)\}$}

\item 如果集合 $A=\{x\mid y=x^2,\,x\in\R\}$，$B=\{y\mid y=x^2,\,x\in\R\}$，则集合 $A$ 与 $B$ 的关系是 \blankline[2em].

% 注：本行答案照原卷作「$B\subset A$」（原卷的 $\subset$ 不带下划线，语义即真包含，与 $\subsetneq$ 等价）；
%     本稿原写作 $\subsetneq$，2026-09-15 回原图核对时按原卷记号改回 $\subset$。
\ans{$B\subset A$}

% 注：本题答案照原卷作「$A\subset B$」（原卷此处的 $\subset$ 不带下划线，即真包含）——
%     $A=\left\{\frac{m}{6}\,\middle|\,m\text{ 为奇数}\right\}$、$B=\left\{\frac{m}{6}\,\middle|\,m\in\Z\right\}$，确是 $A\subsetneq B$。
%     本稿曾误抄成「$A\supsetneq B$（或 $A=B$ 之等价条件）」——方向与结论都反了，
%     2026-09-15 回原图核对时已按原卷订正。
\item 已知 $A=\left\{x\,\middle|\,x=\dfrac{k}{3}-\dfrac{1}{2},\,k\in\Z\right\}$，$B=\left\{x\,\middle|\,x=\dfrac{k}{6}+\dfrac{2}{3},\,k\in\Z\right\}$，则集合 $A$ 与 $B$ 的关系是 \blankline[2em].

\ans{$A\subset B$}

\item 若非空集合 $\{x\mid m\leqslant x\leqslant 2m-\sqrt{m}\}$ 不存在非空真子集，则 $m$ 的取值集合为 \blankline[2em].

\ans{$\{0,1\}$}

\item 已知 $P=\{x\mid 2<x<a,\,x\in\mathbb{N}\}$，且集合 $P$ 中恰有 3 个元素，则实数 $a$ 的取值范围为 \blankline[3em]（用区间表示）.

\ans{$(5,6]$}

\item 设 $x_1,x_2,x_3,x_4$ 是 4 个正整数，从中任取 3 个数求和所得的集合为 $\{25,26,27\}$，则这 4 个数中最小的数为 \blankline[2em].

\ans{$8$}

\ifanswer
\solbegin
\step{从 4 个正整数中任取 3 个数求和后可得到 4 个和，则 4 个和值之和为 $3(x_1+x_2+x_3+x_4)$，必为 3 的倍数.}
\step{又 $27\div 3=9,\,25\div 3=8\dots1,\,26\div 3=8\dots2$，所以这 4 个和为 $25,26,27,27$.}
\step{则 $x_1+x_2+x_3+x_4=\tfrac{1}{3}(25+26+27+27)=35$.}
\step{所以 $35-25=10,\,35-26=9,\,35-27=8$，即这 4 个数分别为 $8,8,9,10$，故这 4 个数中最小的数为 8.}
\solend

\fi
\end{enumerate}

\bigsec{二、选择题}

\begin{enumerate}
\setcounter{enumi}{10}

\item 下列说法中正确的是\mbox{（\quad）.}

\keepwithprev
A. $0$ 与 $\{0\}$ 表示同一个集合；

B. 集合 $\{1,2,3\}$ 与 $\{3,2,1\}$ 是两个相同的集合；

C. 方程 $(x-1)^2(x-2)=0$ 的解集为 $\{1,1,2\}$；

D. 集合 $\{x\mid 4<x<5\}$ 可以用列举法表示.

\ans{B}

\item 下列说法正确的是\mbox{（\quad）.}

\keepwithprev
A. $\varnothing\in\{0\}$\qquad B. $0\subseteq\{0\}$\qquad C. $\varnothing\subseteq\{0,1,2\}$\qquad D. $(1,2)\supset\{1,2\}$

\ans{C}

\item 已知非空集合 $A=\{x\mid x^2-2x+m=0\}$，则集合 $A$ 中所有元素之和为\mbox{（\quad）.}

\keepwithprev
A. $1$\qquad B. $2$\qquad C. $1$ 或 $2$\qquad D. 非以上答案

\ans{C}

\item 如果集合 $S$ 具有性质：① 非空且它的元素都是正整数；② 如果 $x\in S$，那么 $\tfrac{24}{x}\in S$. 请问这样的集合 $S$ 共有\mbox{（\quad）}个.

\keepwithprev
A. $15$ 个\qquad B. $16$ 个\qquad C. $31$ 个\qquad D. $32$ 个

\ans{A}

\end{enumerate}

\bigsec{三、解答题}

\begin{enumerate}
\setcounter{enumi}{14}

\item 设集合 $A=\{x\mid x^2-ax+4=0\}$，$B=\{x\mid x^2-3x+2=0\}$，且 $A\subseteq B$，求实数 $a$ 的取值范围.

\ans{$\{a\mid -4<a\leqslant 4\}$（注：原答案如此）}

\ansspace[6]
\ansneed{7cm}
\item 已知集合 $A=\{x\mid -1<x<6\}$，$B=\{x\mid -m<x<2m+1\}$，其中 $m$ 为实常数.
\begin{enumerate}[label=(\arabic*)]
\item 若 $B\subseteq A$，求实数 $m$ 的取值范围；
% 注：本小问原卷作「求证：$A\subset B$」（不带下划线，即真包含），本稿原写作 $\subseteq$，
%     2026-09-15 回原图核对时按原卷记号订正。
\item 若 $m>3$，求证：$A\subset B$.
\end{enumerate}

\ans{（1）$m\leqslant 1$；（2）略}

\ansspace*[7]
\end{enumerate}

\end{document}
"
%
% 原始材料是微信公众号图文（「魔都数学加油站」，作者署名「破晓」），正文为 3 张 PNG 页
% （1080×1527）；卷面上的粉色斜水印「魔都数学加油站」属水印，未录入。原文本身就是一份
% **讲评版**——题下直接印出【答案】或【解析】，本稿把这些答案与解析移到答案版，试卷版即
% 一张干净的空卷。
%
% 与原卷的排版差异（均已在 README「说明」中交代）：
%   ① 原文自**第 3 题**起（第 1、2 题未在该文发布），源码里以 \setcounter{enumi}{2} 起编；
%   ② 原卷本就印有分节标题「一、填空题」「二、选择题」「三、解答题」（已在原图上逐页放大
%      核对过），本稿照原卷分节，只把标题改排成全稿统一的「大节标题」样式；
%   ③ 原卷的实数集、整数集符号**正体与粗体混用**——第 3 题作正体 $N$、$Z$，第 6 题作粗体
%      $\mathbf{R}$（均逐处放大核对过），本稿按全稿惯例统一写作 $\N$、$\Z$、$\R$；
%   ④ 原卷第 6、7 题的 $\subset$ 不带下划线（即真包含，与带下划线的 $\subseteq$ 在卷面上
%      逐处放大比对后确认不同），本稿按原符号保留，并按真包含口径解答（正文内有 `%` 注）.

\documentclass[a4paper,11pt]{article}
\usepackage{common}

\begin{document}

\examtitlest{2026--2027 学年度位育中学高一上周末作业 1}{集合与常用逻辑用语}

\bigsec{一、填空题}

\begin{enumerate}
\setcounter{enumi}{2}

% 注：本行照原卷录入——原卷作「$y=\dfrac{12}{3-x}\in N,\,x\in Z$」，即「$y=\frac{12}{3-x}$ 是自然数、$x$ 是整数」；
%     本稿曾误写成「$x\in\mathbb{N},\,x\in\mathbb{Z}$」（把 $\in\mathbb{N}$ 挂到了 $x$ 上），
%     那样只能得到 $\{4,6,12\}$，与本卷给的答案 $\{1,2,3,4,6,12\}$ 冲突。
%     2026-09-15 回原图核对时已按原卷订正。
\item 用列举法表示集合 $A=\left\{y\,\middle|\,y=\dfrac{12}{3-x}\in\mathbb{N},\,x\in\mathbb{Z}\right\}=$ \blankline[2em].

\ans{$\{1,2,3,4,6,12\}$}

\item 已知集合 $A=\{2a,\,2a^2+a\}$，若 $1\in A$，则 $a=$ \blankline[2em].

\ans{$-1$.}

\ifanswer
\solbegin
\step{因为 $1\in A$，所以 $2a=1$ 或 $2a^2+a=1$.}
\step{当 $2a=1$ 时，$a=\tfrac{1}{2}$，此时 $2a^2+a=1$，不满足集合中元素的互异性，舍去；}
\step{当 $2a^2+a=1$ 时，解得 $a=-1$ 或 $a=\tfrac{1}{2}$（舍去），此时 $A=\{-2,1\}$.}
\step{综上，$a=-1$.}
\solend

\fi
\item 关于 $x,y$ 的方程组 $\begin{cases} x+y=2 \\ x^2+y^2=2 \end{cases}$ 的解集为 \blankline[2em].

\ans{$\{(1,1)\}$}

\item 如果集合 $A=\{x\mid y=x^2,\,x\in\R\}$，$B=\{y\mid y=x^2,\,x\in\R\}$，则集合 $A$ 与 $B$ 的关系是 \blankline[2em].

% 注：本行答案照原卷作「$B\subset A$」（原卷的 $\subset$ 不带下划线，语义即真包含，与 $\subsetneq$ 等价）；
%     本稿原写作 $\subsetneq$，2026-09-15 回原图核对时按原卷记号改回 $\subset$。
\ans{$B\subset A$}

% 注：本题答案照原卷作「$A\subset B$」（原卷此处的 $\subset$ 不带下划线，即真包含）——
%     $A=\left\{\frac{m}{6}\,\middle|\,m\text{ 为奇数}\right\}$、$B=\left\{\frac{m}{6}\,\middle|\,m\in\Z\right\}$，确是 $A\subsetneq B$。
%     本稿曾误抄成「$A\supsetneq B$（或 $A=B$ 之等价条件）」——方向与结论都反了，
%     2026-09-15 回原图核对时已按原卷订正。
\item 已知 $A=\left\{x\,\middle|\,x=\dfrac{k}{3}-\dfrac{1}{2},\,k\in\Z\right\}$，$B=\left\{x\,\middle|\,x=\dfrac{k}{6}+\dfrac{2}{3},\,k\in\Z\right\}$，则集合 $A$ 与 $B$ 的关系是 \blankline[2em].

\ans{$A\subset B$}

\item 若非空集合 $\{x\mid m\leqslant x\leqslant 2m-\sqrt{m}\}$ 不存在非空真子集，则 $m$ 的取值集合为 \blankline[2em].

\ans{$\{0,1\}$}

\item 已知 $P=\{x\mid 2<x<a,\,x\in\mathbb{N}\}$，且集合 $P$ 中恰有 3 个元素，则实数 $a$ 的取值范围为 \blankline[3em]（用区间表示）.

\ans{$(5,6]$}

\item 设 $x_1,x_2,x_3,x_4$ 是 4 个正整数，从中任取 3 个数求和所得的集合为 $\{25,26,27\}$，则这 4 个数中最小的数为 \blankline[2em].

\ans{$8$}

\ifanswer
\solbegin
\step{从 4 个正整数中任取 3 个数求和后可得到 4 个和，则 4 个和值之和为 $3(x_1+x_2+x_3+x_4)$，必为 3 的倍数.}
\step{又 $27\div 3=9,\,25\div 3=8\dots1,\,26\div 3=8\dots2$，所以这 4 个和为 $25,26,27,27$.}
\step{则 $x_1+x_2+x_3+x_4=\tfrac{1}{3}(25+26+27+27)=35$.}
\step{所以 $35-25=10,\,35-26=9,\,35-27=8$，即这 4 个数分别为 $8,8,9,10$，故这 4 个数中最小的数为 8.}
\solend

\fi
\end{enumerate}

\bigsec{二、选择题}

\begin{enumerate}
\setcounter{enumi}{10}

\item 下列说法中正确的是\mbox{（\quad）.}

\keepwithprev
A. $0$ 与 $\{0\}$ 表示同一个集合；

B. 集合 $\{1,2,3\}$ 与 $\{3,2,1\}$ 是两个相同的集合；

C. 方程 $(x-1)^2(x-2)=0$ 的解集为 $\{1,1,2\}$；

D. 集合 $\{x\mid 4<x<5\}$ 可以用列举法表示.

\ans{B}

\item 下列说法正确的是\mbox{（\quad）.}

\keepwithprev
A. $\varnothing\in\{0\}$\qquad B. $0\subseteq\{0\}$\qquad C. $\varnothing\subseteq\{0,1,2\}$\qquad D. $(1,2)\supset\{1,2\}$

\ans{C}

\item 已知非空集合 $A=\{x\mid x^2-2x+m=0\}$，则集合 $A$ 中所有元素之和为\mbox{（\quad）.}

\keepwithprev
A. $1$\qquad B. $2$\qquad C. $1$ 或 $2$\qquad D. 非以上答案

\ans{C}

\item 如果集合 $S$ 具有性质：① 非空且它的元素都是正整数；② 如果 $x\in S$，那么 $\tfrac{24}{x}\in S$. 请问这样的集合 $S$ 共有\mbox{（\quad）}个.

\keepwithprev
A. $15$ 个\qquad B. $16$ 个\qquad C. $31$ 个\qquad D. $32$ 个

\ans{A}

\end{enumerate}

\bigsec{三、解答题}

\begin{enumerate}
\setcounter{enumi}{14}

\item 设集合 $A=\{x\mid x^2-ax+4=0\}$，$B=\{x\mid x^2-3x+2=0\}$，且 $A\subseteq B$，求实数 $a$ 的取值范围.

\ans{$\{a\mid -4<a\leqslant 4\}$（注：原答案如此）}

\ansspace[6]
\ansneed{7cm}
\item 已知集合 $A=\{x\mid -1<x<6\}$，$B=\{x\mid -m<x<2m+1\}$，其中 $m$ 为实常数.
\begin{enumerate}[label=(\arabic*)]
\item 若 $B\subseteq A$，求实数 $m$ 的取值范围；
% 注：本小问原卷作「求证：$A\subset B$」（不带下划线，即真包含），本稿原写作 $\subseteq$，
%     2026-09-15 回原图核对时按原卷记号订正。
\item 若 $m>3$，求证：$A\subset B$.
\end{enumerate}

\ans{（1）$m\leqslant 1$；（2）略}

\ansspace*[7]
\end{enumerate}

\end{document}
"
%
% 原始材料是微信公众号图文（「魔都数学加油站」，作者署名「破晓」），正文为 3 张 PNG 页
% （1080×1527）；卷面上的粉色斜水印「魔都数学加油站」属水印，未录入。原文本身就是一份
% **讲评版**——题下直接印出【答案】或【解析】，本稿把这些答案与解析移到答案版，试卷版即
% 一张干净的空卷。
%
% 与原卷的排版差异（均已在 README「说明」中交代）：
%   ① 原文自**第 3 题**起（第 1、2 题未在该文发布），源码里以 \setcounter{enumi}{2} 起编；
%   ② 原卷本就印有分节标题「一、填空题」「二、选择题」「三、解答题」（已在原图上逐页放大
%      核对过），本稿照原卷分节，只把标题改排成全稿统一的「大节标题」样式；
%   ③ 原卷的实数集、整数集符号**正体与粗体混用**——第 3 题作正体 $N$、$Z$，第 6 题作粗体
%      $\mathbf{R}$（均逐处放大核对过），本稿按全稿惯例统一写作 $\N$、$\Z$、$\R$；
%   ④ 原卷第 6、7 题的 $\subset$ 不带下划线（即真包含，与带下划线的 $\subseteq$ 在卷面上
%      逐处放大比对后确认不同），本稿按原符号保留，并按真包含口径解答（正文内有 `%` 注）.

\documentclass[a4paper,11pt]{article}
\usepackage{common}

\begin{document}

\examtitlest{2026--2027 学年度位育中学高一上周末作业 1}{集合与常用逻辑用语}

\bigsec{一、填空题}

\begin{enumerate}
\setcounter{enumi}{2}

% 注：本行照原卷录入——原卷作「$y=\dfrac{12}{3-x}\in N,\,x\in Z$」，即「$y=\frac{12}{3-x}$ 是自然数、$x$ 是整数」；
%     本稿曾误写成「$x\in\mathbb{N},\,x\in\mathbb{Z}$」（把 $\in\mathbb{N}$ 挂到了 $x$ 上），
%     那样只能得到 $\{4,6,12\}$，与本卷给的答案 $\{1,2,3,4,6,12\}$ 冲突。
%     2026-09-15 回原图核对时已按原卷订正。
\item 用列举法表示集合 $A=\left\{y\,\middle|\,y=\dfrac{12}{3-x}\in\mathbb{N},\,x\in\mathbb{Z}\right\}=$ \blankline[2em].

\ans{$\{1,2,3,4,6,12\}$}

\item 已知集合 $A=\{2a,\,2a^2+a\}$，若 $1\in A$，则 $a=$ \blankline[2em].

\ans{$-1$.}

\ifanswer
\solbegin
\step{因为 $1\in A$，所以 $2a=1$ 或 $2a^2+a=1$.}
\step{当 $2a=1$ 时，$a=\tfrac{1}{2}$，此时 $2a^2+a=1$，不满足集合中元素的互异性，舍去；}
\step{当 $2a^2+a=1$ 时，解得 $a=-1$ 或 $a=\tfrac{1}{2}$（舍去），此时 $A=\{-2,1\}$.}
\step{综上，$a=-1$.}
\solend

\fi
\item 关于 $x,y$ 的方程组 $\begin{cases} x+y=2 \\ x^2+y^2=2 \end{cases}$ 的解集为 \blankline[2em].

\ans{$\{(1,1)\}$}

\item 如果集合 $A=\{x\mid y=x^2,\,x\in\R\}$，$B=\{y\mid y=x^2,\,x\in\R\}$，则集合 $A$ 与 $B$ 的关系是 \blankline[2em].

% 注：本行答案照原卷作「$B\subset A$」（原卷的 $\subset$ 不带下划线，语义即真包含，与 $\subsetneq$ 等价）；
%     本稿原写作 $\subsetneq$，2026-09-15 回原图核对时按原卷记号改回 $\subset$。
\ans{$B\subset A$}

% 注：本题答案照原卷作「$A\subset B$」（原卷此处的 $\subset$ 不带下划线，即真包含）——
%     $A=\left\{\frac{m}{6}\,\middle|\,m\text{ 为奇数}\right\}$、$B=\left\{\frac{m}{6}\,\middle|\,m\in\Z\right\}$，确是 $A\subsetneq B$。
%     本稿曾误抄成「$A\supsetneq B$（或 $A=B$ 之等价条件）」——方向与结论都反了，
%     2026-09-15 回原图核对时已按原卷订正。
\item 已知 $A=\left\{x\,\middle|\,x=\dfrac{k}{3}-\dfrac{1}{2},\,k\in\Z\right\}$，$B=\left\{x\,\middle|\,x=\dfrac{k}{6}+\dfrac{2}{3},\,k\in\Z\right\}$，则集合 $A$ 与 $B$ 的关系是 \blankline[2em].

\ans{$A\subset B$}

\item 若非空集合 $\{x\mid m\leqslant x\leqslant 2m-\sqrt{m}\}$ 不存在非空真子集，则 $m$ 的取值集合为 \blankline[2em].

\ans{$\{0,1\}$}

\item 已知 $P=\{x\mid 2<x<a,\,x\in\mathbb{N}\}$，且集合 $P$ 中恰有 3 个元素，则实数 $a$ 的取值范围为 \blankline[3em]（用区间表示）.

\ans{$(5,6]$}

\item 设 $x_1,x_2,x_3,x_4$ 是 4 个正整数，从中任取 3 个数求和所得的集合为 $\{25,26,27\}$，则这 4 个数中最小的数为 \blankline[2em].

\ans{$8$}

\ifanswer
\solbegin
\step{从 4 个正整数中任取 3 个数求和后可得到 4 个和，则 4 个和值之和为 $3(x_1+x_2+x_3+x_4)$，必为 3 的倍数.}
\step{又 $27\div 3=9,\,25\div 3=8\dots1,\,26\div 3=8\dots2$，所以这 4 个和为 $25,26,27,27$.}
\step{则 $x_1+x_2+x_3+x_4=\tfrac{1}{3}(25+26+27+27)=35$.}
\step{所以 $35-25=10,\,35-26=9,\,35-27=8$，即这 4 个数分别为 $8,8,9,10$，故这 4 个数中最小的数为 8.}
\solend

\fi
\end{enumerate}

\bigsec{二、选择题}

\begin{enumerate}
\setcounter{enumi}{10}

\item 下列说法中正确的是\mbox{（\quad）.}

\keepwithprev
A. $0$ 与 $\{0\}$ 表示同一个集合；

B. 集合 $\{1,2,3\}$ 与 $\{3,2,1\}$ 是两个相同的集合；

C. 方程 $(x-1)^2(x-2)=0$ 的解集为 $\{1,1,2\}$；

D. 集合 $\{x\mid 4<x<5\}$ 可以用列举法表示.

\ans{B}

\item 下列说法正确的是\mbox{（\quad）.}

\keepwithprev
A. $\varnothing\in\{0\}$\qquad B. $0\subseteq\{0\}$\qquad C. $\varnothing\subseteq\{0,1,2\}$\qquad D. $(1,2)\supset\{1,2\}$

\ans{C}

\item 已知非空集合 $A=\{x\mid x^2-2x+m=0\}$，则集合 $A$ 中所有元素之和为\mbox{（\quad）.}

\keepwithprev
A. $1$\qquad B. $2$\qquad C. $1$ 或 $2$\qquad D. 非以上答案

\ans{C}

\item 如果集合 $S$ 具有性质：① 非空且它的元素都是正整数；② 如果 $x\in S$，那么 $\tfrac{24}{x}\in S$. 请问这样的集合 $S$ 共有\mbox{（\quad）}个.

\keepwithprev
A. $15$ 个\qquad B. $16$ 个\qquad C. $31$ 个\qquad D. $32$ 个

\ans{A}

\end{enumerate}

\bigsec{三、解答题}

\begin{enumerate}
\setcounter{enumi}{14}

\item 设集合 $A=\{x\mid x^2-ax+4=0\}$，$B=\{x\mid x^2-3x+2=0\}$，且 $A\subseteq B$，求实数 $a$ 的取值范围.

\ans{$\{a\mid -4<a\leqslant 4\}$（注：原答案如此）}

\ansspace[6]
\ansneed{7cm}
\item 已知集合 $A=\{x\mid -1<x<6\}$，$B=\{x\mid -m<x<2m+1\}$，其中 $m$ 为实常数.
\begin{enumerate}[label=(\arabic*)]
\item 若 $B\subseteq A$，求实数 $m$ 的取值范围；
% 注：本小问原卷作「求证：$A\subset B$」（不带下划线，即真包含），本稿原写作 $\subseteq$，
%     2026-09-15 回原图核对时按原卷记号订正。
\item 若 $m>3$，求证：$A\subset B$.
\end{enumerate}

\ans{（1）$m\leqslant 1$；（2）略}

\ansspace*[7]
\end{enumerate}

\end{document}
"
% 紧凑版：xelatex "\def\iscompact{1}% 高一04_位育中学_周末作业1.tex —— 高一上数学「2026-2027 年位育中学高一上周末作业 1」（试卷版/答案版）
% 试卷版：xelatex 高一04_位育中学_周末作业1.tex
% 答案版：xelatex "\def\isanswer{1}% 高一04_位育中学_周末作业1.tex —— 高一上数学「2026-2027 年位育中学高一上周末作业 1」（试卷版/答案版）
% 试卷版：xelatex 高一04_位育中学_周末作业1.tex
% 答案版：xelatex "\def\isanswer{1}% 高一04_位育中学_周末作业1.tex —— 高一上数学「2026-2027 年位育中学高一上周末作业 1」（试卷版/答案版）
% 试卷版：xelatex 高一04_位育中学_周末作业1.tex
% 答案版：xelatex "\def\isanswer{1}\input{高一04_位育中学_周末作业1.tex}"
% 紧凑版：xelatex "\def\iscompact{1}\input{高一04_位育中学_周末作业1.tex}"
%
% 原始材料是微信公众号图文（「魔都数学加油站」，作者署名「破晓」），正文为 3 张 PNG 页
% （1080×1527）；卷面上的粉色斜水印「魔都数学加油站」属水印，未录入。原文本身就是一份
% **讲评版**——题下直接印出【答案】或【解析】，本稿把这些答案与解析移到答案版，试卷版即
% 一张干净的空卷。
%
% 与原卷的排版差异（均已在 README「说明」中交代）：
%   ① 原文自**第 3 题**起（第 1、2 题未在该文发布），源码里以 \setcounter{enumi}{2} 起编；
%   ② 原卷本就印有分节标题「一、填空题」「二、选择题」「三、解答题」（已在原图上逐页放大
%      核对过），本稿照原卷分节，只把标题改排成全稿统一的「大节标题」样式；
%   ③ 原卷的实数集、整数集符号**正体与粗体混用**——第 3 题作正体 $N$、$Z$，第 6 题作粗体
%      $\mathbf{R}$（均逐处放大核对过），本稿按全稿惯例统一写作 $\N$、$\Z$、$\R$；
%   ④ 原卷第 6、7 题的 $\subset$ 不带下划线（即真包含，与带下划线的 $\subseteq$ 在卷面上
%      逐处放大比对后确认不同），本稿按原符号保留，并按真包含口径解答（正文内有 `%` 注）.

\documentclass[a4paper,11pt]{article}
\usepackage{common}

\begin{document}

\examtitlest{2026--2027 学年度位育中学高一上周末作业 1}{集合与常用逻辑用语}

\bigsec{一、填空题}

\begin{enumerate}
\setcounter{enumi}{2}

% 注：本行照原卷录入——原卷作「$y=\dfrac{12}{3-x}\in N,\,x\in Z$」，即「$y=\frac{12}{3-x}$ 是自然数、$x$ 是整数」；
%     本稿曾误写成「$x\in\mathbb{N},\,x\in\mathbb{Z}$」（把 $\in\mathbb{N}$ 挂到了 $x$ 上），
%     那样只能得到 $\{4,6,12\}$，与本卷给的答案 $\{1,2,3,4,6,12\}$ 冲突。
%     2026-09-15 回原图核对时已按原卷订正。
\item 用列举法表示集合 $A=\left\{y\,\middle|\,y=\dfrac{12}{3-x}\in\mathbb{N},\,x\in\mathbb{Z}\right\}=$ \blankline[2em].

\ans{$\{1,2,3,4,6,12\}$}

\item 已知集合 $A=\{2a,\,2a^2+a\}$，若 $1\in A$，则 $a=$ \blankline[2em].

\ans{$-1$.}

\ifanswer
\solbegin
\step{因为 $1\in A$，所以 $2a=1$ 或 $2a^2+a=1$.}
\step{当 $2a=1$ 时，$a=\tfrac{1}{2}$，此时 $2a^2+a=1$，不满足集合中元素的互异性，舍去；}
\step{当 $2a^2+a=1$ 时，解得 $a=-1$ 或 $a=\tfrac{1}{2}$（舍去），此时 $A=\{-2,1\}$.}
\step{综上，$a=-1$.}
\solend

\fi
\item 关于 $x,y$ 的方程组 $\begin{cases} x+y=2 \\ x^2+y^2=2 \end{cases}$ 的解集为 \blankline[2em].

\ans{$\{(1,1)\}$}

\item 如果集合 $A=\{x\mid y=x^2,\,x\in\R\}$，$B=\{y\mid y=x^2,\,x\in\R\}$，则集合 $A$ 与 $B$ 的关系是 \blankline[2em].

% 注：本行答案照原卷作「$B\subset A$」（原卷的 $\subset$ 不带下划线，语义即真包含，与 $\subsetneq$ 等价）；
%     本稿原写作 $\subsetneq$，2026-09-15 回原图核对时按原卷记号改回 $\subset$。
\ans{$B\subset A$}

% 注：本题答案照原卷作「$A\subset B$」（原卷此处的 $\subset$ 不带下划线，即真包含）——
%     $A=\left\{\frac{m}{6}\,\middle|\,m\text{ 为奇数}\right\}$、$B=\left\{\frac{m}{6}\,\middle|\,m\in\Z\right\}$，确是 $A\subsetneq B$。
%     本稿曾误抄成「$A\supsetneq B$（或 $A=B$ 之等价条件）」——方向与结论都反了，
%     2026-09-15 回原图核对时已按原卷订正。
\item 已知 $A=\left\{x\,\middle|\,x=\dfrac{k}{3}-\dfrac{1}{2},\,k\in\Z\right\}$，$B=\left\{x\,\middle|\,x=\dfrac{k}{6}+\dfrac{2}{3},\,k\in\Z\right\}$，则集合 $A$ 与 $B$ 的关系是 \blankline[2em].

\ans{$A\subset B$}

\item 若非空集合 $\{x\mid m\leqslant x\leqslant 2m-\sqrt{m}\}$ 不存在非空真子集，则 $m$ 的取值集合为 \blankline[2em].

\ans{$\{0,1\}$}

\item 已知 $P=\{x\mid 2<x<a,\,x\in\mathbb{N}\}$，且集合 $P$ 中恰有 3 个元素，则实数 $a$ 的取值范围为 \blankline[3em]（用区间表示）.

\ans{$(5,6]$}

\item 设 $x_1,x_2,x_3,x_4$ 是 4 个正整数，从中任取 3 个数求和所得的集合为 $\{25,26,27\}$，则这 4 个数中最小的数为 \blankline[2em].

\ans{$8$}

\ifanswer
\solbegin
\step{从 4 个正整数中任取 3 个数求和后可得到 4 个和，则 4 个和值之和为 $3(x_1+x_2+x_3+x_4)$，必为 3 的倍数.}
\step{又 $27\div 3=9,\,25\div 3=8\dots1,\,26\div 3=8\dots2$，所以这 4 个和为 $25,26,27,27$.}
\step{则 $x_1+x_2+x_3+x_4=\tfrac{1}{3}(25+26+27+27)=35$.}
\step{所以 $35-25=10,\,35-26=9,\,35-27=8$，即这 4 个数分别为 $8,8,9,10$，故这 4 个数中最小的数为 8.}
\solend

\fi
\end{enumerate}

\bigsec{二、选择题}

\begin{enumerate}
\setcounter{enumi}{10}

\item 下列说法中正确的是\mbox{（\quad）.}

\keepwithprev
A. $0$ 与 $\{0\}$ 表示同一个集合；

B. 集合 $\{1,2,3\}$ 与 $\{3,2,1\}$ 是两个相同的集合；

C. 方程 $(x-1)^2(x-2)=0$ 的解集为 $\{1,1,2\}$；

D. 集合 $\{x\mid 4<x<5\}$ 可以用列举法表示.

\ans{B}

\item 下列说法正确的是\mbox{（\quad）.}

\keepwithprev
A. $\varnothing\in\{0\}$\qquad B. $0\subseteq\{0\}$\qquad C. $\varnothing\subseteq\{0,1,2\}$\qquad D. $(1,2)\supset\{1,2\}$

\ans{C}

\item 已知非空集合 $A=\{x\mid x^2-2x+m=0\}$，则集合 $A$ 中所有元素之和为\mbox{（\quad）.}

\keepwithprev
A. $1$\qquad B. $2$\qquad C. $1$ 或 $2$\qquad D. 非以上答案

\ans{C}

\item 如果集合 $S$ 具有性质：① 非空且它的元素都是正整数；② 如果 $x\in S$，那么 $\tfrac{24}{x}\in S$. 请问这样的集合 $S$ 共有\mbox{（\quad）}个.

\keepwithprev
A. $15$ 个\qquad B. $16$ 个\qquad C. $31$ 个\qquad D. $32$ 个

\ans{A}

\end{enumerate}

\bigsec{三、解答题}

\begin{enumerate}
\setcounter{enumi}{14}

\item 设集合 $A=\{x\mid x^2-ax+4=0\}$，$B=\{x\mid x^2-3x+2=0\}$，且 $A\subseteq B$，求实数 $a$ 的取值范围.

\ans{$\{a\mid -4<a\leqslant 4\}$（注：原答案如此）}

\ansspace[6]
\ansneed{7cm}
\item 已知集合 $A=\{x\mid -1<x<6\}$，$B=\{x\mid -m<x<2m+1\}$，其中 $m$ 为实常数.
\begin{enumerate}[label=(\arabic*)]
\item 若 $B\subseteq A$，求实数 $m$ 的取值范围；
% 注：本小问原卷作「求证：$A\subset B$」（不带下划线，即真包含），本稿原写作 $\subseteq$，
%     2026-09-15 回原图核对时按原卷记号订正。
\item 若 $m>3$，求证：$A\subset B$.
\end{enumerate}

\ans{（1）$m\leqslant 1$；（2）略}

\ansspace*[7]
\end{enumerate}

\end{document}
"
% 紧凑版：xelatex "\def\iscompact{1}% 高一04_位育中学_周末作业1.tex —— 高一上数学「2026-2027 年位育中学高一上周末作业 1」（试卷版/答案版）
% 试卷版：xelatex 高一04_位育中学_周末作业1.tex
% 答案版：xelatex "\def\isanswer{1}\input{高一04_位育中学_周末作业1.tex}"
% 紧凑版：xelatex "\def\iscompact{1}\input{高一04_位育中学_周末作业1.tex}"
%
% 原始材料是微信公众号图文（「魔都数学加油站」，作者署名「破晓」），正文为 3 张 PNG 页
% （1080×1527）；卷面上的粉色斜水印「魔都数学加油站」属水印，未录入。原文本身就是一份
% **讲评版**——题下直接印出【答案】或【解析】，本稿把这些答案与解析移到答案版，试卷版即
% 一张干净的空卷。
%
% 与原卷的排版差异（均已在 README「说明」中交代）：
%   ① 原文自**第 3 题**起（第 1、2 题未在该文发布），源码里以 \setcounter{enumi}{2} 起编；
%   ② 原卷本就印有分节标题「一、填空题」「二、选择题」「三、解答题」（已在原图上逐页放大
%      核对过），本稿照原卷分节，只把标题改排成全稿统一的「大节标题」样式；
%   ③ 原卷的实数集、整数集符号**正体与粗体混用**——第 3 题作正体 $N$、$Z$，第 6 题作粗体
%      $\mathbf{R}$（均逐处放大核对过），本稿按全稿惯例统一写作 $\N$、$\Z$、$\R$；
%   ④ 原卷第 6、7 题的 $\subset$ 不带下划线（即真包含，与带下划线的 $\subseteq$ 在卷面上
%      逐处放大比对后确认不同），本稿按原符号保留，并按真包含口径解答（正文内有 `%` 注）.

\documentclass[a4paper,11pt]{article}
\usepackage{common}

\begin{document}

\examtitlest{2026--2027 学年度位育中学高一上周末作业 1}{集合与常用逻辑用语}

\bigsec{一、填空题}

\begin{enumerate}
\setcounter{enumi}{2}

% 注：本行照原卷录入——原卷作「$y=\dfrac{12}{3-x}\in N,\,x\in Z$」，即「$y=\frac{12}{3-x}$ 是自然数、$x$ 是整数」；
%     本稿曾误写成「$x\in\mathbb{N},\,x\in\mathbb{Z}$」（把 $\in\mathbb{N}$ 挂到了 $x$ 上），
%     那样只能得到 $\{4,6,12\}$，与本卷给的答案 $\{1,2,3,4,6,12\}$ 冲突。
%     2026-09-15 回原图核对时已按原卷订正。
\item 用列举法表示集合 $A=\left\{y\,\middle|\,y=\dfrac{12}{3-x}\in\mathbb{N},\,x\in\mathbb{Z}\right\}=$ \blankline[2em].

\ans{$\{1,2,3,4,6,12\}$}

\item 已知集合 $A=\{2a,\,2a^2+a\}$，若 $1\in A$，则 $a=$ \blankline[2em].

\ans{$-1$.}

\ifanswer
\solbegin
\step{因为 $1\in A$，所以 $2a=1$ 或 $2a^2+a=1$.}
\step{当 $2a=1$ 时，$a=\tfrac{1}{2}$，此时 $2a^2+a=1$，不满足集合中元素的互异性，舍去；}
\step{当 $2a^2+a=1$ 时，解得 $a=-1$ 或 $a=\tfrac{1}{2}$（舍去），此时 $A=\{-2,1\}$.}
\step{综上，$a=-1$.}
\solend

\fi
\item 关于 $x,y$ 的方程组 $\begin{cases} x+y=2 \\ x^2+y^2=2 \end{cases}$ 的解集为 \blankline[2em].

\ans{$\{(1,1)\}$}

\item 如果集合 $A=\{x\mid y=x^2,\,x\in\R\}$，$B=\{y\mid y=x^2,\,x\in\R\}$，则集合 $A$ 与 $B$ 的关系是 \blankline[2em].

% 注：本行答案照原卷作「$B\subset A$」（原卷的 $\subset$ 不带下划线，语义即真包含，与 $\subsetneq$ 等价）；
%     本稿原写作 $\subsetneq$，2026-09-15 回原图核对时按原卷记号改回 $\subset$。
\ans{$B\subset A$}

% 注：本题答案照原卷作「$A\subset B$」（原卷此处的 $\subset$ 不带下划线，即真包含）——
%     $A=\left\{\frac{m}{6}\,\middle|\,m\text{ 为奇数}\right\}$、$B=\left\{\frac{m}{6}\,\middle|\,m\in\Z\right\}$，确是 $A\subsetneq B$。
%     本稿曾误抄成「$A\supsetneq B$（或 $A=B$ 之等价条件）」——方向与结论都反了，
%     2026-09-15 回原图核对时已按原卷订正。
\item 已知 $A=\left\{x\,\middle|\,x=\dfrac{k}{3}-\dfrac{1}{2},\,k\in\Z\right\}$，$B=\left\{x\,\middle|\,x=\dfrac{k}{6}+\dfrac{2}{3},\,k\in\Z\right\}$，则集合 $A$ 与 $B$ 的关系是 \blankline[2em].

\ans{$A\subset B$}

\item 若非空集合 $\{x\mid m\leqslant x\leqslant 2m-\sqrt{m}\}$ 不存在非空真子集，则 $m$ 的取值集合为 \blankline[2em].

\ans{$\{0,1\}$}

\item 已知 $P=\{x\mid 2<x<a,\,x\in\mathbb{N}\}$，且集合 $P$ 中恰有 3 个元素，则实数 $a$ 的取值范围为 \blankline[3em]（用区间表示）.

\ans{$(5,6]$}

\item 设 $x_1,x_2,x_3,x_4$ 是 4 个正整数，从中任取 3 个数求和所得的集合为 $\{25,26,27\}$，则这 4 个数中最小的数为 \blankline[2em].

\ans{$8$}

\ifanswer
\solbegin
\step{从 4 个正整数中任取 3 个数求和后可得到 4 个和，则 4 个和值之和为 $3(x_1+x_2+x_3+x_4)$，必为 3 的倍数.}
\step{又 $27\div 3=9,\,25\div 3=8\dots1,\,26\div 3=8\dots2$，所以这 4 个和为 $25,26,27,27$.}
\step{则 $x_1+x_2+x_3+x_4=\tfrac{1}{3}(25+26+27+27)=35$.}
\step{所以 $35-25=10,\,35-26=9,\,35-27=8$，即这 4 个数分别为 $8,8,9,10$，故这 4 个数中最小的数为 8.}
\solend

\fi
\end{enumerate}

\bigsec{二、选择题}

\begin{enumerate}
\setcounter{enumi}{10}

\item 下列说法中正确的是\mbox{（\quad）.}

\keepwithprev
A. $0$ 与 $\{0\}$ 表示同一个集合；

B. 集合 $\{1,2,3\}$ 与 $\{3,2,1\}$ 是两个相同的集合；

C. 方程 $(x-1)^2(x-2)=0$ 的解集为 $\{1,1,2\}$；

D. 集合 $\{x\mid 4<x<5\}$ 可以用列举法表示.

\ans{B}

\item 下列说法正确的是\mbox{（\quad）.}

\keepwithprev
A. $\varnothing\in\{0\}$\qquad B. $0\subseteq\{0\}$\qquad C. $\varnothing\subseteq\{0,1,2\}$\qquad D. $(1,2)\supset\{1,2\}$

\ans{C}

\item 已知非空集合 $A=\{x\mid x^2-2x+m=0\}$，则集合 $A$ 中所有元素之和为\mbox{（\quad）.}

\keepwithprev
A. $1$\qquad B. $2$\qquad C. $1$ 或 $2$\qquad D. 非以上答案

\ans{C}

\item 如果集合 $S$ 具有性质：① 非空且它的元素都是正整数；② 如果 $x\in S$，那么 $\tfrac{24}{x}\in S$. 请问这样的集合 $S$ 共有\mbox{（\quad）}个.

\keepwithprev
A. $15$ 个\qquad B. $16$ 个\qquad C. $31$ 个\qquad D. $32$ 个

\ans{A}

\end{enumerate}

\bigsec{三、解答题}

\begin{enumerate}
\setcounter{enumi}{14}

\item 设集合 $A=\{x\mid x^2-ax+4=0\}$，$B=\{x\mid x^2-3x+2=0\}$，且 $A\subseteq B$，求实数 $a$ 的取值范围.

\ans{$\{a\mid -4<a\leqslant 4\}$（注：原答案如此）}

\ansspace[6]
\ansneed{7cm}
\item 已知集合 $A=\{x\mid -1<x<6\}$，$B=\{x\mid -m<x<2m+1\}$，其中 $m$ 为实常数.
\begin{enumerate}[label=(\arabic*)]
\item 若 $B\subseteq A$，求实数 $m$ 的取值范围；
% 注：本小问原卷作「求证：$A\subset B$」（不带下划线，即真包含），本稿原写作 $\subseteq$，
%     2026-09-15 回原图核对时按原卷记号订正。
\item 若 $m>3$，求证：$A\subset B$.
\end{enumerate}

\ans{（1）$m\leqslant 1$；（2）略}

\ansspace*[7]
\end{enumerate}

\end{document}
"
%
% 原始材料是微信公众号图文（「魔都数学加油站」，作者署名「破晓」），正文为 3 张 PNG 页
% （1080×1527）；卷面上的粉色斜水印「魔都数学加油站」属水印，未录入。原文本身就是一份
% **讲评版**——题下直接印出【答案】或【解析】，本稿把这些答案与解析移到答案版，试卷版即
% 一张干净的空卷。
%
% 与原卷的排版差异（均已在 README「说明」中交代）：
%   ① 原文自**第 3 题**起（第 1、2 题未在该文发布），源码里以 \setcounter{enumi}{2} 起编；
%   ② 原卷本就印有分节标题「一、填空题」「二、选择题」「三、解答题」（已在原图上逐页放大
%      核对过），本稿照原卷分节，只把标题改排成全稿统一的「大节标题」样式；
%   ③ 原卷的实数集、整数集符号**正体与粗体混用**——第 3 题作正体 $N$、$Z$，第 6 题作粗体
%      $\mathbf{R}$（均逐处放大核对过），本稿按全稿惯例统一写作 $\N$、$\Z$、$\R$；
%   ④ 原卷第 6、7 题的 $\subset$ 不带下划线（即真包含，与带下划线的 $\subseteq$ 在卷面上
%      逐处放大比对后确认不同），本稿按原符号保留，并按真包含口径解答（正文内有 `%` 注）.

\documentclass[a4paper,11pt]{article}
\usepackage{common}

\begin{document}

\examtitlest{2026--2027 学年度位育中学高一上周末作业 1}{集合与常用逻辑用语}

\bigsec{一、填空题}

\begin{enumerate}
\setcounter{enumi}{2}

% 注：本行照原卷录入——原卷作「$y=\dfrac{12}{3-x}\in N,\,x\in Z$」，即「$y=\frac{12}{3-x}$ 是自然数、$x$ 是整数」；
%     本稿曾误写成「$x\in\mathbb{N},\,x\in\mathbb{Z}$」（把 $\in\mathbb{N}$ 挂到了 $x$ 上），
%     那样只能得到 $\{4,6,12\}$，与本卷给的答案 $\{1,2,3,4,6,12\}$ 冲突。
%     2026-09-15 回原图核对时已按原卷订正。
\item 用列举法表示集合 $A=\left\{y\,\middle|\,y=\dfrac{12}{3-x}\in\mathbb{N},\,x\in\mathbb{Z}\right\}=$ \blankline[2em].

\ans{$\{1,2,3,4,6,12\}$}

\item 已知集合 $A=\{2a,\,2a^2+a\}$，若 $1\in A$，则 $a=$ \blankline[2em].

\ans{$-1$.}

\ifanswer
\solbegin
\step{因为 $1\in A$，所以 $2a=1$ 或 $2a^2+a=1$.}
\step{当 $2a=1$ 时，$a=\tfrac{1}{2}$，此时 $2a^2+a=1$，不满足集合中元素的互异性，舍去；}
\step{当 $2a^2+a=1$ 时，解得 $a=-1$ 或 $a=\tfrac{1}{2}$（舍去），此时 $A=\{-2,1\}$.}
\step{综上，$a=-1$.}
\solend

\fi
\item 关于 $x,y$ 的方程组 $\begin{cases} x+y=2 \\ x^2+y^2=2 \end{cases}$ 的解集为 \blankline[2em].

\ans{$\{(1,1)\}$}

\item 如果集合 $A=\{x\mid y=x^2,\,x\in\R\}$，$B=\{y\mid y=x^2,\,x\in\R\}$，则集合 $A$ 与 $B$ 的关系是 \blankline[2em].

% 注：本行答案照原卷作「$B\subset A$」（原卷的 $\subset$ 不带下划线，语义即真包含，与 $\subsetneq$ 等价）；
%     本稿原写作 $\subsetneq$，2026-09-15 回原图核对时按原卷记号改回 $\subset$。
\ans{$B\subset A$}

% 注：本题答案照原卷作「$A\subset B$」（原卷此处的 $\subset$ 不带下划线，即真包含）——
%     $A=\left\{\frac{m}{6}\,\middle|\,m\text{ 为奇数}\right\}$、$B=\left\{\frac{m}{6}\,\middle|\,m\in\Z\right\}$，确是 $A\subsetneq B$。
%     本稿曾误抄成「$A\supsetneq B$（或 $A=B$ 之等价条件）」——方向与结论都反了，
%     2026-09-15 回原图核对时已按原卷订正。
\item 已知 $A=\left\{x\,\middle|\,x=\dfrac{k}{3}-\dfrac{1}{2},\,k\in\Z\right\}$，$B=\left\{x\,\middle|\,x=\dfrac{k}{6}+\dfrac{2}{3},\,k\in\Z\right\}$，则集合 $A$ 与 $B$ 的关系是 \blankline[2em].

\ans{$A\subset B$}

\item 若非空集合 $\{x\mid m\leqslant x\leqslant 2m-\sqrt{m}\}$ 不存在非空真子集，则 $m$ 的取值集合为 \blankline[2em].

\ans{$\{0,1\}$}

\item 已知 $P=\{x\mid 2<x<a,\,x\in\mathbb{N}\}$，且集合 $P$ 中恰有 3 个元素，则实数 $a$ 的取值范围为 \blankline[3em]（用区间表示）.

\ans{$(5,6]$}

\item 设 $x_1,x_2,x_3,x_4$ 是 4 个正整数，从中任取 3 个数求和所得的集合为 $\{25,26,27\}$，则这 4 个数中最小的数为 \blankline[2em].

\ans{$8$}

\ifanswer
\solbegin
\step{从 4 个正整数中任取 3 个数求和后可得到 4 个和，则 4 个和值之和为 $3(x_1+x_2+x_3+x_4)$，必为 3 的倍数.}
\step{又 $27\div 3=9,\,25\div 3=8\dots1,\,26\div 3=8\dots2$，所以这 4 个和为 $25,26,27,27$.}
\step{则 $x_1+x_2+x_3+x_4=\tfrac{1}{3}(25+26+27+27)=35$.}
\step{所以 $35-25=10,\,35-26=9,\,35-27=8$，即这 4 个数分别为 $8,8,9,10$，故这 4 个数中最小的数为 8.}
\solend

\fi
\end{enumerate}

\bigsec{二、选择题}

\begin{enumerate}
\setcounter{enumi}{10}

\item 下列说法中正确的是\mbox{（\quad）.}

\keepwithprev
A. $0$ 与 $\{0\}$ 表示同一个集合；

B. 集合 $\{1,2,3\}$ 与 $\{3,2,1\}$ 是两个相同的集合；

C. 方程 $(x-1)^2(x-2)=0$ 的解集为 $\{1,1,2\}$；

D. 集合 $\{x\mid 4<x<5\}$ 可以用列举法表示.

\ans{B}

\item 下列说法正确的是\mbox{（\quad）.}

\keepwithprev
A. $\varnothing\in\{0\}$\qquad B. $0\subseteq\{0\}$\qquad C. $\varnothing\subseteq\{0,1,2\}$\qquad D. $(1,2)\supset\{1,2\}$

\ans{C}

\item 已知非空集合 $A=\{x\mid x^2-2x+m=0\}$，则集合 $A$ 中所有元素之和为\mbox{（\quad）.}

\keepwithprev
A. $1$\qquad B. $2$\qquad C. $1$ 或 $2$\qquad D. 非以上答案

\ans{C}

\item 如果集合 $S$ 具有性质：① 非空且它的元素都是正整数；② 如果 $x\in S$，那么 $\tfrac{24}{x}\in S$. 请问这样的集合 $S$ 共有\mbox{（\quad）}个.

\keepwithprev
A. $15$ 个\qquad B. $16$ 个\qquad C. $31$ 个\qquad D. $32$ 个

\ans{A}

\end{enumerate}

\bigsec{三、解答题}

\begin{enumerate}
\setcounter{enumi}{14}

\item 设集合 $A=\{x\mid x^2-ax+4=0\}$，$B=\{x\mid x^2-3x+2=0\}$，且 $A\subseteq B$，求实数 $a$ 的取值范围.

\ans{$\{a\mid -4<a\leqslant 4\}$（注：原答案如此）}

\ansspace[6]
\ansneed{7cm}
\item 已知集合 $A=\{x\mid -1<x<6\}$，$B=\{x\mid -m<x<2m+1\}$，其中 $m$ 为实常数.
\begin{enumerate}[label=(\arabic*)]
\item 若 $B\subseteq A$，求实数 $m$ 的取值范围；
% 注：本小问原卷作「求证：$A\subset B$」（不带下划线，即真包含），本稿原写作 $\subseteq$，
%     2026-09-15 回原图核对时按原卷记号订正。
\item 若 $m>3$，求证：$A\subset B$.
\end{enumerate}

\ans{（1）$m\leqslant 1$；（2）略}

\ansspace*[7]
\end{enumerate}

\end{document}
"
% 紧凑版：xelatex "\def\iscompact{1}% 高一04_位育中学_周末作业1.tex —— 高一上数学「2026-2027 年位育中学高一上周末作业 1」（试卷版/答案版）
% 试卷版：xelatex 高一04_位育中学_周末作业1.tex
% 答案版：xelatex "\def\isanswer{1}% 高一04_位育中学_周末作业1.tex —— 高一上数学「2026-2027 年位育中学高一上周末作业 1」（试卷版/答案版）
% 试卷版：xelatex 高一04_位育中学_周末作业1.tex
% 答案版：xelatex "\def\isanswer{1}\input{高一04_位育中学_周末作业1.tex}"
% 紧凑版：xelatex "\def\iscompact{1}\input{高一04_位育中学_周末作业1.tex}"
%
% 原始材料是微信公众号图文（「魔都数学加油站」，作者署名「破晓」），正文为 3 张 PNG 页
% （1080×1527）；卷面上的粉色斜水印「魔都数学加油站」属水印，未录入。原文本身就是一份
% **讲评版**——题下直接印出【答案】或【解析】，本稿把这些答案与解析移到答案版，试卷版即
% 一张干净的空卷。
%
% 与原卷的排版差异（均已在 README「说明」中交代）：
%   ① 原文自**第 3 题**起（第 1、2 题未在该文发布），源码里以 \setcounter{enumi}{2} 起编；
%   ② 原卷本就印有分节标题「一、填空题」「二、选择题」「三、解答题」（已在原图上逐页放大
%      核对过），本稿照原卷分节，只把标题改排成全稿统一的「大节标题」样式；
%   ③ 原卷的实数集、整数集符号**正体与粗体混用**——第 3 题作正体 $N$、$Z$，第 6 题作粗体
%      $\mathbf{R}$（均逐处放大核对过），本稿按全稿惯例统一写作 $\N$、$\Z$、$\R$；
%   ④ 原卷第 6、7 题的 $\subset$ 不带下划线（即真包含，与带下划线的 $\subseteq$ 在卷面上
%      逐处放大比对后确认不同），本稿按原符号保留，并按真包含口径解答（正文内有 `%` 注）.

\documentclass[a4paper,11pt]{article}
\usepackage{common}

\begin{document}

\examtitlest{2026--2027 学年度位育中学高一上周末作业 1}{集合与常用逻辑用语}

\bigsec{一、填空题}

\begin{enumerate}
\setcounter{enumi}{2}

% 注：本行照原卷录入——原卷作「$y=\dfrac{12}{3-x}\in N,\,x\in Z$」，即「$y=\frac{12}{3-x}$ 是自然数、$x$ 是整数」；
%     本稿曾误写成「$x\in\mathbb{N},\,x\in\mathbb{Z}$」（把 $\in\mathbb{N}$ 挂到了 $x$ 上），
%     那样只能得到 $\{4,6,12\}$，与本卷给的答案 $\{1,2,3,4,6,12\}$ 冲突。
%     2026-09-15 回原图核对时已按原卷订正。
\item 用列举法表示集合 $A=\left\{y\,\middle|\,y=\dfrac{12}{3-x}\in\mathbb{N},\,x\in\mathbb{Z}\right\}=$ \blankline[2em].

\ans{$\{1,2,3,4,6,12\}$}

\item 已知集合 $A=\{2a,\,2a^2+a\}$，若 $1\in A$，则 $a=$ \blankline[2em].

\ans{$-1$.}

\ifanswer
\solbegin
\step{因为 $1\in A$，所以 $2a=1$ 或 $2a^2+a=1$.}
\step{当 $2a=1$ 时，$a=\tfrac{1}{2}$，此时 $2a^2+a=1$，不满足集合中元素的互异性，舍去；}
\step{当 $2a^2+a=1$ 时，解得 $a=-1$ 或 $a=\tfrac{1}{2}$（舍去），此时 $A=\{-2,1\}$.}
\step{综上，$a=-1$.}
\solend

\fi
\item 关于 $x,y$ 的方程组 $\begin{cases} x+y=2 \\ x^2+y^2=2 \end{cases}$ 的解集为 \blankline[2em].

\ans{$\{(1,1)\}$}

\item 如果集合 $A=\{x\mid y=x^2,\,x\in\R\}$，$B=\{y\mid y=x^2,\,x\in\R\}$，则集合 $A$ 与 $B$ 的关系是 \blankline[2em].

% 注：本行答案照原卷作「$B\subset A$」（原卷的 $\subset$ 不带下划线，语义即真包含，与 $\subsetneq$ 等价）；
%     本稿原写作 $\subsetneq$，2026-09-15 回原图核对时按原卷记号改回 $\subset$。
\ans{$B\subset A$}

% 注：本题答案照原卷作「$A\subset B$」（原卷此处的 $\subset$ 不带下划线，即真包含）——
%     $A=\left\{\frac{m}{6}\,\middle|\,m\text{ 为奇数}\right\}$、$B=\left\{\frac{m}{6}\,\middle|\,m\in\Z\right\}$，确是 $A\subsetneq B$。
%     本稿曾误抄成「$A\supsetneq B$（或 $A=B$ 之等价条件）」——方向与结论都反了，
%     2026-09-15 回原图核对时已按原卷订正。
\item 已知 $A=\left\{x\,\middle|\,x=\dfrac{k}{3}-\dfrac{1}{2},\,k\in\Z\right\}$，$B=\left\{x\,\middle|\,x=\dfrac{k}{6}+\dfrac{2}{3},\,k\in\Z\right\}$，则集合 $A$ 与 $B$ 的关系是 \blankline[2em].

\ans{$A\subset B$}

\item 若非空集合 $\{x\mid m\leqslant x\leqslant 2m-\sqrt{m}\}$ 不存在非空真子集，则 $m$ 的取值集合为 \blankline[2em].

\ans{$\{0,1\}$}

\item 已知 $P=\{x\mid 2<x<a,\,x\in\mathbb{N}\}$，且集合 $P$ 中恰有 3 个元素，则实数 $a$ 的取值范围为 \blankline[3em]（用区间表示）.

\ans{$(5,6]$}

\item 设 $x_1,x_2,x_3,x_4$ 是 4 个正整数，从中任取 3 个数求和所得的集合为 $\{25,26,27\}$，则这 4 个数中最小的数为 \blankline[2em].

\ans{$8$}

\ifanswer
\solbegin
\step{从 4 个正整数中任取 3 个数求和后可得到 4 个和，则 4 个和值之和为 $3(x_1+x_2+x_3+x_4)$，必为 3 的倍数.}
\step{又 $27\div 3=9,\,25\div 3=8\dots1,\,26\div 3=8\dots2$，所以这 4 个和为 $25,26,27,27$.}
\step{则 $x_1+x_2+x_3+x_4=\tfrac{1}{3}(25+26+27+27)=35$.}
\step{所以 $35-25=10,\,35-26=9,\,35-27=8$，即这 4 个数分别为 $8,8,9,10$，故这 4 个数中最小的数为 8.}
\solend

\fi
\end{enumerate}

\bigsec{二、选择题}

\begin{enumerate}
\setcounter{enumi}{10}

\item 下列说法中正确的是\mbox{（\quad）.}

\keepwithprev
A. $0$ 与 $\{0\}$ 表示同一个集合；

B. 集合 $\{1,2,3\}$ 与 $\{3,2,1\}$ 是两个相同的集合；

C. 方程 $(x-1)^2(x-2)=0$ 的解集为 $\{1,1,2\}$；

D. 集合 $\{x\mid 4<x<5\}$ 可以用列举法表示.

\ans{B}

\item 下列说法正确的是\mbox{（\quad）.}

\keepwithprev
A. $\varnothing\in\{0\}$\qquad B. $0\subseteq\{0\}$\qquad C. $\varnothing\subseteq\{0,1,2\}$\qquad D. $(1,2)\supset\{1,2\}$

\ans{C}

\item 已知非空集合 $A=\{x\mid x^2-2x+m=0\}$，则集合 $A$ 中所有元素之和为\mbox{（\quad）.}

\keepwithprev
A. $1$\qquad B. $2$\qquad C. $1$ 或 $2$\qquad D. 非以上答案

\ans{C}

\item 如果集合 $S$ 具有性质：① 非空且它的元素都是正整数；② 如果 $x\in S$，那么 $\tfrac{24}{x}\in S$. 请问这样的集合 $S$ 共有\mbox{（\quad）}个.

\keepwithprev
A. $15$ 个\qquad B. $16$ 个\qquad C. $31$ 个\qquad D. $32$ 个

\ans{A}

\end{enumerate}

\bigsec{三、解答题}

\begin{enumerate}
\setcounter{enumi}{14}

\item 设集合 $A=\{x\mid x^2-ax+4=0\}$，$B=\{x\mid x^2-3x+2=0\}$，且 $A\subseteq B$，求实数 $a$ 的取值范围.

\ans{$\{a\mid -4<a\leqslant 4\}$（注：原答案如此）}

\ansspace[6]
\ansneed{7cm}
\item 已知集合 $A=\{x\mid -1<x<6\}$，$B=\{x\mid -m<x<2m+1\}$，其中 $m$ 为实常数.
\begin{enumerate}[label=(\arabic*)]
\item 若 $B\subseteq A$，求实数 $m$ 的取值范围；
% 注：本小问原卷作「求证：$A\subset B$」（不带下划线，即真包含），本稿原写作 $\subseteq$，
%     2026-09-15 回原图核对时按原卷记号订正。
\item 若 $m>3$，求证：$A\subset B$.
\end{enumerate}

\ans{（1）$m\leqslant 1$；（2）略}

\ansspace*[7]
\end{enumerate}

\end{document}
"
% 紧凑版：xelatex "\def\iscompact{1}% 高一04_位育中学_周末作业1.tex —— 高一上数学「2026-2027 年位育中学高一上周末作业 1」（试卷版/答案版）
% 试卷版：xelatex 高一04_位育中学_周末作业1.tex
% 答案版：xelatex "\def\isanswer{1}\input{高一04_位育中学_周末作业1.tex}"
% 紧凑版：xelatex "\def\iscompact{1}\input{高一04_位育中学_周末作业1.tex}"
%
% 原始材料是微信公众号图文（「魔都数学加油站」，作者署名「破晓」），正文为 3 张 PNG 页
% （1080×1527）；卷面上的粉色斜水印「魔都数学加油站」属水印，未录入。原文本身就是一份
% **讲评版**——题下直接印出【答案】或【解析】，本稿把这些答案与解析移到答案版，试卷版即
% 一张干净的空卷。
%
% 与原卷的排版差异（均已在 README「说明」中交代）：
%   ① 原文自**第 3 题**起（第 1、2 题未在该文发布），源码里以 \setcounter{enumi}{2} 起编；
%   ② 原卷本就印有分节标题「一、填空题」「二、选择题」「三、解答题」（已在原图上逐页放大
%      核对过），本稿照原卷分节，只把标题改排成全稿统一的「大节标题」样式；
%   ③ 原卷的实数集、整数集符号**正体与粗体混用**——第 3 题作正体 $N$、$Z$，第 6 题作粗体
%      $\mathbf{R}$（均逐处放大核对过），本稿按全稿惯例统一写作 $\N$、$\Z$、$\R$；
%   ④ 原卷第 6、7 题的 $\subset$ 不带下划线（即真包含，与带下划线的 $\subseteq$ 在卷面上
%      逐处放大比对后确认不同），本稿按原符号保留，并按真包含口径解答（正文内有 `%` 注）.

\documentclass[a4paper,11pt]{article}
\usepackage{common}

\begin{document}

\examtitlest{2026--2027 学年度位育中学高一上周末作业 1}{集合与常用逻辑用语}

\bigsec{一、填空题}

\begin{enumerate}
\setcounter{enumi}{2}

% 注：本行照原卷录入——原卷作「$y=\dfrac{12}{3-x}\in N,\,x\in Z$」，即「$y=\frac{12}{3-x}$ 是自然数、$x$ 是整数」；
%     本稿曾误写成「$x\in\mathbb{N},\,x\in\mathbb{Z}$」（把 $\in\mathbb{N}$ 挂到了 $x$ 上），
%     那样只能得到 $\{4,6,12\}$，与本卷给的答案 $\{1,2,3,4,6,12\}$ 冲突。
%     2026-09-15 回原图核对时已按原卷订正。
\item 用列举法表示集合 $A=\left\{y\,\middle|\,y=\dfrac{12}{3-x}\in\mathbb{N},\,x\in\mathbb{Z}\right\}=$ \blankline[2em].

\ans{$\{1,2,3,4,6,12\}$}

\item 已知集合 $A=\{2a,\,2a^2+a\}$，若 $1\in A$，则 $a=$ \blankline[2em].

\ans{$-1$.}

\ifanswer
\solbegin
\step{因为 $1\in A$，所以 $2a=1$ 或 $2a^2+a=1$.}
\step{当 $2a=1$ 时，$a=\tfrac{1}{2}$，此时 $2a^2+a=1$，不满足集合中元素的互异性，舍去；}
\step{当 $2a^2+a=1$ 时，解得 $a=-1$ 或 $a=\tfrac{1}{2}$（舍去），此时 $A=\{-2,1\}$.}
\step{综上，$a=-1$.}
\solend

\fi
\item 关于 $x,y$ 的方程组 $\begin{cases} x+y=2 \\ x^2+y^2=2 \end{cases}$ 的解集为 \blankline[2em].

\ans{$\{(1,1)\}$}

\item 如果集合 $A=\{x\mid y=x^2,\,x\in\R\}$，$B=\{y\mid y=x^2,\,x\in\R\}$，则集合 $A$ 与 $B$ 的关系是 \blankline[2em].

% 注：本行答案照原卷作「$B\subset A$」（原卷的 $\subset$ 不带下划线，语义即真包含，与 $\subsetneq$ 等价）；
%     本稿原写作 $\subsetneq$，2026-09-15 回原图核对时按原卷记号改回 $\subset$。
\ans{$B\subset A$}

% 注：本题答案照原卷作「$A\subset B$」（原卷此处的 $\subset$ 不带下划线，即真包含）——
%     $A=\left\{\frac{m}{6}\,\middle|\,m\text{ 为奇数}\right\}$、$B=\left\{\frac{m}{6}\,\middle|\,m\in\Z\right\}$，确是 $A\subsetneq B$。
%     本稿曾误抄成「$A\supsetneq B$（或 $A=B$ 之等价条件）」——方向与结论都反了，
%     2026-09-15 回原图核对时已按原卷订正。
\item 已知 $A=\left\{x\,\middle|\,x=\dfrac{k}{3}-\dfrac{1}{2},\,k\in\Z\right\}$，$B=\left\{x\,\middle|\,x=\dfrac{k}{6}+\dfrac{2}{3},\,k\in\Z\right\}$，则集合 $A$ 与 $B$ 的关系是 \blankline[2em].

\ans{$A\subset B$}

\item 若非空集合 $\{x\mid m\leqslant x\leqslant 2m-\sqrt{m}\}$ 不存在非空真子集，则 $m$ 的取值集合为 \blankline[2em].

\ans{$\{0,1\}$}

\item 已知 $P=\{x\mid 2<x<a,\,x\in\mathbb{N}\}$，且集合 $P$ 中恰有 3 个元素，则实数 $a$ 的取值范围为 \blankline[3em]（用区间表示）.

\ans{$(5,6]$}

\item 设 $x_1,x_2,x_3,x_4$ 是 4 个正整数，从中任取 3 个数求和所得的集合为 $\{25,26,27\}$，则这 4 个数中最小的数为 \blankline[2em].

\ans{$8$}

\ifanswer
\solbegin
\step{从 4 个正整数中任取 3 个数求和后可得到 4 个和，则 4 个和值之和为 $3(x_1+x_2+x_3+x_4)$，必为 3 的倍数.}
\step{又 $27\div 3=9,\,25\div 3=8\dots1,\,26\div 3=8\dots2$，所以这 4 个和为 $25,26,27,27$.}
\step{则 $x_1+x_2+x_3+x_4=\tfrac{1}{3}(25+26+27+27)=35$.}
\step{所以 $35-25=10,\,35-26=9,\,35-27=8$，即这 4 个数分别为 $8,8,9,10$，故这 4 个数中最小的数为 8.}
\solend

\fi
\end{enumerate}

\bigsec{二、选择题}

\begin{enumerate}
\setcounter{enumi}{10}

\item 下列说法中正确的是\mbox{（\quad）.}

\keepwithprev
A. $0$ 与 $\{0\}$ 表示同一个集合；

B. 集合 $\{1,2,3\}$ 与 $\{3,2,1\}$ 是两个相同的集合；

C. 方程 $(x-1)^2(x-2)=0$ 的解集为 $\{1,1,2\}$；

D. 集合 $\{x\mid 4<x<5\}$ 可以用列举法表示.

\ans{B}

\item 下列说法正确的是\mbox{（\quad）.}

\keepwithprev
A. $\varnothing\in\{0\}$\qquad B. $0\subseteq\{0\}$\qquad C. $\varnothing\subseteq\{0,1,2\}$\qquad D. $(1,2)\supset\{1,2\}$

\ans{C}

\item 已知非空集合 $A=\{x\mid x^2-2x+m=0\}$，则集合 $A$ 中所有元素之和为\mbox{（\quad）.}

\keepwithprev
A. $1$\qquad B. $2$\qquad C. $1$ 或 $2$\qquad D. 非以上答案

\ans{C}

\item 如果集合 $S$ 具有性质：① 非空且它的元素都是正整数；② 如果 $x\in S$，那么 $\tfrac{24}{x}\in S$. 请问这样的集合 $S$ 共有\mbox{（\quad）}个.

\keepwithprev
A. $15$ 个\qquad B. $16$ 个\qquad C. $31$ 个\qquad D. $32$ 个

\ans{A}

\end{enumerate}

\bigsec{三、解答题}

\begin{enumerate}
\setcounter{enumi}{14}

\item 设集合 $A=\{x\mid x^2-ax+4=0\}$，$B=\{x\mid x^2-3x+2=0\}$，且 $A\subseteq B$，求实数 $a$ 的取值范围.

\ans{$\{a\mid -4<a\leqslant 4\}$（注：原答案如此）}

\ansspace[6]
\ansneed{7cm}
\item 已知集合 $A=\{x\mid -1<x<6\}$，$B=\{x\mid -m<x<2m+1\}$，其中 $m$ 为实常数.
\begin{enumerate}[label=(\arabic*)]
\item 若 $B\subseteq A$，求实数 $m$ 的取值范围；
% 注：本小问原卷作「求证：$A\subset B$」（不带下划线，即真包含），本稿原写作 $\subseteq$，
%     2026-09-15 回原图核对时按原卷记号订正。
\item 若 $m>3$，求证：$A\subset B$.
\end{enumerate}

\ans{（1）$m\leqslant 1$；（2）略}

\ansspace*[7]
\end{enumerate}

\end{document}
"
%
% 原始材料是微信公众号图文（「魔都数学加油站」，作者署名「破晓」），正文为 3 张 PNG 页
% （1080×1527）；卷面上的粉色斜水印「魔都数学加油站」属水印，未录入。原文本身就是一份
% **讲评版**——题下直接印出【答案】或【解析】，本稿把这些答案与解析移到答案版，试卷版即
% 一张干净的空卷。
%
% 与原卷的排版差异（均已在 README「说明」中交代）：
%   ① 原文自**第 3 题**起（第 1、2 题未在该文发布），源码里以 \setcounter{enumi}{2} 起编；
%   ② 原卷本就印有分节标题「一、填空题」「二、选择题」「三、解答题」（已在原图上逐页放大
%      核对过），本稿照原卷分节，只把标题改排成全稿统一的「大节标题」样式；
%   ③ 原卷的实数集、整数集符号**正体与粗体混用**——第 3 题作正体 $N$、$Z$，第 6 题作粗体
%      $\mathbf{R}$（均逐处放大核对过），本稿按全稿惯例统一写作 $\N$、$\Z$、$\R$；
%   ④ 原卷第 6、7 题的 $\subset$ 不带下划线（即真包含，与带下划线的 $\subseteq$ 在卷面上
%      逐处放大比对后确认不同），本稿按原符号保留，并按真包含口径解答（正文内有 `%` 注）.

\documentclass[a4paper,11pt]{article}
\usepackage{common}

\begin{document}

\examtitlest{2026--2027 学年度位育中学高一上周末作业 1}{集合与常用逻辑用语}

\bigsec{一、填空题}

\begin{enumerate}
\setcounter{enumi}{2}

% 注：本行照原卷录入——原卷作「$y=\dfrac{12}{3-x}\in N,\,x\in Z$」，即「$y=\frac{12}{3-x}$ 是自然数、$x$ 是整数」；
%     本稿曾误写成「$x\in\mathbb{N},\,x\in\mathbb{Z}$」（把 $\in\mathbb{N}$ 挂到了 $x$ 上），
%     那样只能得到 $\{4,6,12\}$，与本卷给的答案 $\{1,2,3,4,6,12\}$ 冲突。
%     2026-09-15 回原图核对时已按原卷订正。
\item 用列举法表示集合 $A=\left\{y\,\middle|\,y=\dfrac{12}{3-x}\in\mathbb{N},\,x\in\mathbb{Z}\right\}=$ \blankline[2em].

\ans{$\{1,2,3,4,6,12\}$}

\item 已知集合 $A=\{2a,\,2a^2+a\}$，若 $1\in A$，则 $a=$ \blankline[2em].

\ans{$-1$.}

\ifanswer
\solbegin
\step{因为 $1\in A$，所以 $2a=1$ 或 $2a^2+a=1$.}
\step{当 $2a=1$ 时，$a=\tfrac{1}{2}$，此时 $2a^2+a=1$，不满足集合中元素的互异性，舍去；}
\step{当 $2a^2+a=1$ 时，解得 $a=-1$ 或 $a=\tfrac{1}{2}$（舍去），此时 $A=\{-2,1\}$.}
\step{综上，$a=-1$.}
\solend

\fi
\item 关于 $x,y$ 的方程组 $\begin{cases} x+y=2 \\ x^2+y^2=2 \end{cases}$ 的解集为 \blankline[2em].

\ans{$\{(1,1)\}$}

\item 如果集合 $A=\{x\mid y=x^2,\,x\in\R\}$，$B=\{y\mid y=x^2,\,x\in\R\}$，则集合 $A$ 与 $B$ 的关系是 \blankline[2em].

% 注：本行答案照原卷作「$B\subset A$」（原卷的 $\subset$ 不带下划线，语义即真包含，与 $\subsetneq$ 等价）；
%     本稿原写作 $\subsetneq$，2026-09-15 回原图核对时按原卷记号改回 $\subset$。
\ans{$B\subset A$}

% 注：本题答案照原卷作「$A\subset B$」（原卷此处的 $\subset$ 不带下划线，即真包含）——
%     $A=\left\{\frac{m}{6}\,\middle|\,m\text{ 为奇数}\right\}$、$B=\left\{\frac{m}{6}\,\middle|\,m\in\Z\right\}$，确是 $A\subsetneq B$。
%     本稿曾误抄成「$A\supsetneq B$（或 $A=B$ 之等价条件）」——方向与结论都反了，
%     2026-09-15 回原图核对时已按原卷订正。
\item 已知 $A=\left\{x\,\middle|\,x=\dfrac{k}{3}-\dfrac{1}{2},\,k\in\Z\right\}$，$B=\left\{x\,\middle|\,x=\dfrac{k}{6}+\dfrac{2}{3},\,k\in\Z\right\}$，则集合 $A$ 与 $B$ 的关系是 \blankline[2em].

\ans{$A\subset B$}

\item 若非空集合 $\{x\mid m\leqslant x\leqslant 2m-\sqrt{m}\}$ 不存在非空真子集，则 $m$ 的取值集合为 \blankline[2em].

\ans{$\{0,1\}$}

\item 已知 $P=\{x\mid 2<x<a,\,x\in\mathbb{N}\}$，且集合 $P$ 中恰有 3 个元素，则实数 $a$ 的取值范围为 \blankline[3em]（用区间表示）.

\ans{$(5,6]$}

\item 设 $x_1,x_2,x_3,x_4$ 是 4 个正整数，从中任取 3 个数求和所得的集合为 $\{25,26,27\}$，则这 4 个数中最小的数为 \blankline[2em].

\ans{$8$}

\ifanswer
\solbegin
\step{从 4 个正整数中任取 3 个数求和后可得到 4 个和，则 4 个和值之和为 $3(x_1+x_2+x_3+x_4)$，必为 3 的倍数.}
\step{又 $27\div 3=9,\,25\div 3=8\dots1,\,26\div 3=8\dots2$，所以这 4 个和为 $25,26,27,27$.}
\step{则 $x_1+x_2+x_3+x_4=\tfrac{1}{3}(25+26+27+27)=35$.}
\step{所以 $35-25=10,\,35-26=9,\,35-27=8$，即这 4 个数分别为 $8,8,9,10$，故这 4 个数中最小的数为 8.}
\solend

\fi
\end{enumerate}

\bigsec{二、选择题}

\begin{enumerate}
\setcounter{enumi}{10}

\item 下列说法中正确的是\mbox{（\quad）.}

\keepwithprev
A. $0$ 与 $\{0\}$ 表示同一个集合；

B. 集合 $\{1,2,3\}$ 与 $\{3,2,1\}$ 是两个相同的集合；

C. 方程 $(x-1)^2(x-2)=0$ 的解集为 $\{1,1,2\}$；

D. 集合 $\{x\mid 4<x<5\}$ 可以用列举法表示.

\ans{B}

\item 下列说法正确的是\mbox{（\quad）.}

\keepwithprev
A. $\varnothing\in\{0\}$\qquad B. $0\subseteq\{0\}$\qquad C. $\varnothing\subseteq\{0,1,2\}$\qquad D. $(1,2)\supset\{1,2\}$

\ans{C}

\item 已知非空集合 $A=\{x\mid x^2-2x+m=0\}$，则集合 $A$ 中所有元素之和为\mbox{（\quad）.}

\keepwithprev
A. $1$\qquad B. $2$\qquad C. $1$ 或 $2$\qquad D. 非以上答案

\ans{C}

\item 如果集合 $S$ 具有性质：① 非空且它的元素都是正整数；② 如果 $x\in S$，那么 $\tfrac{24}{x}\in S$. 请问这样的集合 $S$ 共有\mbox{（\quad）}个.

\keepwithprev
A. $15$ 个\qquad B. $16$ 个\qquad C. $31$ 个\qquad D. $32$ 个

\ans{A}

\end{enumerate}

\bigsec{三、解答题}

\begin{enumerate}
\setcounter{enumi}{14}

\item 设集合 $A=\{x\mid x^2-ax+4=0\}$，$B=\{x\mid x^2-3x+2=0\}$，且 $A\subseteq B$，求实数 $a$ 的取值范围.

\ans{$\{a\mid -4<a\leqslant 4\}$（注：原答案如此）}

\ansspace[6]
\ansneed{7cm}
\item 已知集合 $A=\{x\mid -1<x<6\}$，$B=\{x\mid -m<x<2m+1\}$，其中 $m$ 为实常数.
\begin{enumerate}[label=(\arabic*)]
\item 若 $B\subseteq A$，求实数 $m$ 的取值范围；
% 注：本小问原卷作「求证：$A\subset B$」（不带下划线，即真包含），本稿原写作 $\subseteq$，
%     2026-09-15 回原图核对时按原卷记号订正。
\item 若 $m>3$，求证：$A\subset B$.
\end{enumerate}

\ans{（1）$m\leqslant 1$；（2）略}

\ansspace*[7]
\end{enumerate}

\end{document}
"
%
% 原始材料是微信公众号图文（「魔都数学加油站」，作者署名「破晓」），正文为 3 张 PNG 页
% （1080×1527）；卷面上的粉色斜水印「魔都数学加油站」属水印，未录入。原文本身就是一份
% **讲评版**——题下直接印出【答案】或【解析】，本稿把这些答案与解析移到答案版，试卷版即
% 一张干净的空卷。
%
% 与原卷的排版差异（均已在 README「说明」中交代）：
%   ① 原文自**第 3 题**起（第 1、2 题未在该文发布），源码里以 \setcounter{enumi}{2} 起编；
%   ② 原卷本就印有分节标题「一、填空题」「二、选择题」「三、解答题」（已在原图上逐页放大
%      核对过），本稿照原卷分节，只把标题改排成全稿统一的「大节标题」样式；
%   ③ 原卷的实数集、整数集符号**正体与粗体混用**——第 3 题作正体 $N$、$Z$，第 6 题作粗体
%      $\mathbf{R}$（均逐处放大核对过），本稿按全稿惯例统一写作 $\N$、$\Z$、$\R$；
%   ④ 原卷第 6、7 题的 $\subset$ 不带下划线（即真包含，与带下划线的 $\subseteq$ 在卷面上
%      逐处放大比对后确认不同），本稿按原符号保留，并按真包含口径解答（正文内有 `%` 注）.

\documentclass[a4paper,11pt]{article}
\usepackage{common}

\begin{document}

\examtitlest{2026--2027 学年度位育中学高一上周末作业 1}{集合与常用逻辑用语}

\bigsec{一、填空题}

\begin{enumerate}
\setcounter{enumi}{2}

% 注：本行照原卷录入——原卷作「$y=\dfrac{12}{3-x}\in N,\,x\in Z$」，即「$y=\frac{12}{3-x}$ 是自然数、$x$ 是整数」；
%     本稿曾误写成「$x\in\mathbb{N},\,x\in\mathbb{Z}$」（把 $\in\mathbb{N}$ 挂到了 $x$ 上），
%     那样只能得到 $\{4,6,12\}$，与本卷给的答案 $\{1,2,3,4,6,12\}$ 冲突。
%     2026-09-15 回原图核对时已按原卷订正。
\item 用列举法表示集合 $A=\left\{y\,\middle|\,y=\dfrac{12}{3-x}\in\mathbb{N},\,x\in\mathbb{Z}\right\}=$ \blankline[2em].

\ans{$\{1,2,3,4,6,12\}$}

\item 已知集合 $A=\{2a,\,2a^2+a\}$，若 $1\in A$，则 $a=$ \blankline[2em].

\ans{$-1$.}

\ifanswer
\solbegin
\step{因为 $1\in A$，所以 $2a=1$ 或 $2a^2+a=1$.}
\step{当 $2a=1$ 时，$a=\tfrac{1}{2}$，此时 $2a^2+a=1$，不满足集合中元素的互异性，舍去；}
\step{当 $2a^2+a=1$ 时，解得 $a=-1$ 或 $a=\tfrac{1}{2}$（舍去），此时 $A=\{-2,1\}$.}
\step{综上，$a=-1$.}
\solend

\fi
\item 关于 $x,y$ 的方程组 $\begin{cases} x+y=2 \\ x^2+y^2=2 \end{cases}$ 的解集为 \blankline[2em].

\ans{$\{(1,1)\}$}

\item 如果集合 $A=\{x\mid y=x^2,\,x\in\R\}$，$B=\{y\mid y=x^2,\,x\in\R\}$，则集合 $A$ 与 $B$ 的关系是 \blankline[2em].

% 注：本行答案照原卷作「$B\subset A$」（原卷的 $\subset$ 不带下划线，语义即真包含，与 $\subsetneq$ 等价）；
%     本稿原写作 $\subsetneq$，2026-09-15 回原图核对时按原卷记号改回 $\subset$。
\ans{$B\subset A$}

% 注：本题答案照原卷作「$A\subset B$」（原卷此处的 $\subset$ 不带下划线，即真包含）——
%     $A=\left\{\frac{m}{6}\,\middle|\,m\text{ 为奇数}\right\}$、$B=\left\{\frac{m}{6}\,\middle|\,m\in\Z\right\}$，确是 $A\subsetneq B$。
%     本稿曾误抄成「$A\supsetneq B$（或 $A=B$ 之等价条件）」——方向与结论都反了，
%     2026-09-15 回原图核对时已按原卷订正。
\item 已知 $A=\left\{x\,\middle|\,x=\dfrac{k}{3}-\dfrac{1}{2},\,k\in\Z\right\}$，$B=\left\{x\,\middle|\,x=\dfrac{k}{6}+\dfrac{2}{3},\,k\in\Z\right\}$，则集合 $A$ 与 $B$ 的关系是 \blankline[2em].

\ans{$A\subset B$}

\item 若非空集合 $\{x\mid m\leqslant x\leqslant 2m-\sqrt{m}\}$ 不存在非空真子集，则 $m$ 的取值集合为 \blankline[2em].

\ans{$\{0,1\}$}

\item 已知 $P=\{x\mid 2<x<a,\,x\in\mathbb{N}\}$，且集合 $P$ 中恰有 3 个元素，则实数 $a$ 的取值范围为 \blankline[3em]（用区间表示）.

\ans{$(5,6]$}

\item 设 $x_1,x_2,x_3,x_4$ 是 4 个正整数，从中任取 3 个数求和所得的集合为 $\{25,26,27\}$，则这 4 个数中最小的数为 \blankline[2em].

\ans{$8$}

\ifanswer
\solbegin
\step{从 4 个正整数中任取 3 个数求和后可得到 4 个和，则 4 个和值之和为 $3(x_1+x_2+x_3+x_4)$，必为 3 的倍数.}
\step{又 $27\div 3=9,\,25\div 3=8\dots1,\,26\div 3=8\dots2$，所以这 4 个和为 $25,26,27,27$.}
\step{则 $x_1+x_2+x_3+x_4=\tfrac{1}{3}(25+26+27+27)=35$.}
\step{所以 $35-25=10,\,35-26=9,\,35-27=8$，即这 4 个数分别为 $8,8,9,10$，故这 4 个数中最小的数为 8.}
\solend

\fi
\end{enumerate}

\bigsec{二、选择题}

\begin{enumerate}
\setcounter{enumi}{10}

\item 下列说法中正确的是\mbox{（\quad）.}

\keepwithprev
A. $0$ 与 $\{0\}$ 表示同一个集合；

B. 集合 $\{1,2,3\}$ 与 $\{3,2,1\}$ 是两个相同的集合；

C. 方程 $(x-1)^2(x-2)=0$ 的解集为 $\{1,1,2\}$；

D. 集合 $\{x\mid 4<x<5\}$ 可以用列举法表示.

\ans{B}

\item 下列说法正确的是\mbox{（\quad）.}

\keepwithprev
A. $\varnothing\in\{0\}$\qquad B. $0\subseteq\{0\}$\qquad C. $\varnothing\subseteq\{0,1,2\}$\qquad D. $(1,2)\supset\{1,2\}$

\ans{C}

\item 已知非空集合 $A=\{x\mid x^2-2x+m=0\}$，则集合 $A$ 中所有元素之和为\mbox{（\quad）.}

\keepwithprev
A. $1$\qquad B. $2$\qquad C. $1$ 或 $2$\qquad D. 非以上答案

\ans{C}

\item 如果集合 $S$ 具有性质：① 非空且它的元素都是正整数；② 如果 $x\in S$，那么 $\tfrac{24}{x}\in S$. 请问这样的集合 $S$ 共有\mbox{（\quad）}个.

\keepwithprev
A. $15$ 个\qquad B. $16$ 个\qquad C. $31$ 个\qquad D. $32$ 个

\ans{A}

\end{enumerate}

\bigsec{三、解答题}

\begin{enumerate}
\setcounter{enumi}{14}

\item 设集合 $A=\{x\mid x^2-ax+4=0\}$，$B=\{x\mid x^2-3x+2=0\}$，且 $A\subseteq B$，求实数 $a$ 的取值范围.

\ans{$\{a\mid -4<a\leqslant 4\}$（注：原答案如此）}

\ansspace[6]
\ansneed{7cm}
\item 已知集合 $A=\{x\mid -1<x<6\}$，$B=\{x\mid -m<x<2m+1\}$，其中 $m$ 为实常数.
\begin{enumerate}[label=(\arabic*)]
\item 若 $B\subseteq A$，求实数 $m$ 的取值范围；
% 注：本小问原卷作「求证：$A\subset B$」（不带下划线，即真包含），本稿原写作 $\subseteq$，
%     2026-09-15 回原图核对时按原卷记号订正。
\item 若 $m>3$，求证：$A\subset B$.
\end{enumerate}

\ans{（1）$m\leqslant 1$；（2）略}

\ansspace*[7]
\end{enumerate}

\end{document}
"
%
% 原始材料是微信公众号图文（「魔都数学加油站」，作者署名「破晓」），正文为 3 张 PNG 页
% （1080×1527）；卷面上的粉色斜水印「魔都数学加油站」属水印，未录入。原文本身就是一份
% **讲评版**——题下直接印出【答案】或【解析】，本稿把这些答案与解析移到答案版，试卷版即
% 一张干净的空卷。
%
% 与原卷的排版差异（均已在 README「说明」中交代）：
%   ① 原文自**第 3 题**起（第 1、2 题未在该文发布），源码里以 \setcounter{enumi}{2} 起编；
%   ② 原卷本就印有分节标题「一、填空题」「二、选择题」「三、解答题」（已在原图上逐页放大
%      核对过），本稿照原卷分节，只把标题改排成全稿统一的「大节标题」样式；
%   ③ 原卷的实数集、整数集符号**正体与粗体混用**——第 3 题作正体 $N$、$Z$，第 6 题作粗体
%      $\mathbf{R}$（均逐处放大核对过），本稿按全稿惯例统一写作 $\N$、$\Z$、$\R$；
%   ④ 原卷第 6、7 题的 $\subset$ 不带下划线（即真包含，与带下划线的 $\subseteq$ 在卷面上
%      逐处放大比对后确认不同），本稿按原符号保留，并按真包含口径解答（正文内有 `%` 注）.

\documentclass[a4paper,11pt]{article}
\usepackage{common}

\begin{document}

\examtitlest{2026--2027 学年度位育中学高一上周末作业 1}{集合与常用逻辑用语}

\bigsec{一、填空题}

\begin{enumerate}
\setcounter{enumi}{2}

% 注：本行照原卷录入——原卷作「$y=\dfrac{12}{3-x}\in N,\,x\in Z$」，即「$y=\frac{12}{3-x}$ 是自然数、$x$ 是整数」；
%     本稿曾误写成「$x\in\mathbb{N},\,x\in\mathbb{Z}$」（把 $\in\mathbb{N}$ 挂到了 $x$ 上），
%     那样只能得到 $\{4,6,12\}$，与本卷给的答案 $\{1,2,3,4,6,12\}$ 冲突。
%     2026-09-15 回原图核对时已按原卷订正。
\item 用列举法表示集合 $A=\left\{y\,\middle|\,y=\dfrac{12}{3-x}\in\mathbb{N},\,x\in\mathbb{Z}\right\}=$ \blankline[2em].

\ans{$\{1,2,3,4,6,12\}$}

\item 已知集合 $A=\{2a,\,2a^2+a\}$，若 $1\in A$，则 $a=$ \blankline[2em].

\ans{$-1$.}

\ifanswer
\solbegin
\step{因为 $1\in A$，所以 $2a=1$ 或 $2a^2+a=1$.}
\step{当 $2a=1$ 时，$a=\tfrac{1}{2}$，此时 $2a^2+a=1$，不满足集合中元素的互异性，舍去；}
\step{当 $2a^2+a=1$ 时，解得 $a=-1$ 或 $a=\tfrac{1}{2}$（舍去），此时 $A=\{-2,1\}$.}
\step{综上，$a=-1$.}
\solend

\fi
\item 关于 $x,y$ 的方程组 $\begin{cases} x+y=2 \\ x^2+y^2=2 \end{cases}$ 的解集为 \blankline[2em].

\ans{$\{(1,1)\}$}

\item 如果集合 $A=\{x\mid y=x^2,\,x\in\R\}$，$B=\{y\mid y=x^2,\,x\in\R\}$，则集合 $A$ 与 $B$ 的关系是 \blankline[2em].

% 注：本行答案照原卷作「$B\subset A$」（原卷的 $\subset$ 不带下划线，语义即真包含，与 $\subsetneq$ 等价）；
%     本稿原写作 $\subsetneq$，2026-09-15 回原图核对时按原卷记号改回 $\subset$。
\ans{$B\subset A$}

% 注：本题答案照原卷作「$A\subset B$」（原卷此处的 $\subset$ 不带下划线，即真包含）——
%     $A=\left\{\frac{m}{6}\,\middle|\,m\text{ 为奇数}\right\}$、$B=\left\{\frac{m}{6}\,\middle|\,m\in\Z\right\}$，确是 $A\subsetneq B$。
%     本稿曾误抄成「$A\supsetneq B$（或 $A=B$ 之等价条件）」——方向与结论都反了，
%     2026-09-15 回原图核对时已按原卷订正。
\item 已知 $A=\left\{x\,\middle|\,x=\dfrac{k}{3}-\dfrac{1}{2},\,k\in\Z\right\}$，$B=\left\{x\,\middle|\,x=\dfrac{k}{6}+\dfrac{2}{3},\,k\in\Z\right\}$，则集合 $A$ 与 $B$ 的关系是 \blankline[2em].

\ans{$A\subset B$}

\item 若非空集合 $\{x\mid m\leqslant x\leqslant 2m-\sqrt{m}\}$ 不存在非空真子集，则 $m$ 的取值集合为 \blankline[2em].

\ans{$\{0,1\}$}

\item 已知 $P=\{x\mid 2<x<a,\,x\in\mathbb{N}\}$，且集合 $P$ 中恰有 3 个元素，则实数 $a$ 的取值范围为 \blankline[3em]（用区间表示）.

\ans{$(5,6]$}

\item 设 $x_1,x_2,x_3,x_4$ 是 4 个正整数，从中任取 3 个数求和所得的集合为 $\{25,26,27\}$，则这 4 个数中最小的数为 \blankline[2em].

\ans{$8$}

\ifanswer
\solbegin
\step{从 4 个正整数中任取 3 个数求和后可得到 4 个和，则 4 个和值之和为 $3(x_1+x_2+x_3+x_4)$，必为 3 的倍数.}
\step{又 $27\div 3=9,\,25\div 3=8\dots1,\,26\div 3=8\dots2$，所以这 4 个和为 $25,26,27,27$.}
\step{则 $x_1+x_2+x_3+x_4=\tfrac{1}{3}(25+26+27+27)=35$.}
\step{所以 $35-25=10,\,35-26=9,\,35-27=8$，即这 4 个数分别为 $8,8,9,10$，故这 4 个数中最小的数为 8.}
\solend

\fi
\end{enumerate}

\bigsec{二、选择题}

\begin{enumerate}
\setcounter{enumi}{10}

\item 下列说法中正确的是\mbox{（\quad）.}

\keepwithprev
A. $0$ 与 $\{0\}$ 表示同一个集合；

B. 集合 $\{1,2,3\}$ 与 $\{3,2,1\}$ 是两个相同的集合；

C. 方程 $(x-1)^2(x-2)=0$ 的解集为 $\{1,1,2\}$；

D. 集合 $\{x\mid 4<x<5\}$ 可以用列举法表示.

\ans{B}

\item 下列说法正确的是\mbox{（\quad）.}

\keepwithprev
A. $\varnothing\in\{0\}$\qquad B. $0\subseteq\{0\}$\qquad C. $\varnothing\subseteq\{0,1,2\}$\qquad D. $(1,2)\supset\{1,2\}$

\ans{C}

\item 已知非空集合 $A=\{x\mid x^2-2x+m=0\}$，则集合 $A$ 中所有元素之和为\mbox{（\quad）.}

\keepwithprev
A. $1$\qquad B. $2$\qquad C. $1$ 或 $2$\qquad D. 非以上答案

\ans{C}

\item 如果集合 $S$ 具有性质：① 非空且它的元素都是正整数；② 如果 $x\in S$，那么 $\tfrac{24}{x}\in S$. 请问这样的集合 $S$ 共有\mbox{（\quad）}个.

\keepwithprev
A. $15$ 个\qquad B. $16$ 个\qquad C. $31$ 个\qquad D. $32$ 个

\ans{A}

\end{enumerate}

\bigsec{三、解答题}

\begin{enumerate}
\setcounter{enumi}{14}

\item 设集合 $A=\{x\mid x^2-ax+4=0\}$，$B=\{x\mid x^2-3x+2=0\}$，且 $A\subseteq B$，求实数 $a$ 的取值范围.

\ans{$\{a\mid -4<a\leqslant 4\}$（注：原答案如此）}

\ansspace[6]
\ansneed{7cm}
\item 已知集合 $A=\{x\mid -1<x<6\}$，$B=\{x\mid -m<x<2m+1\}$，其中 $m$ 为实常数.
\begin{enumerate}[label=(\arabic*)]
\item 若 $B\subseteq A$，求实数 $m$ 的取值范围；
% 注：本小问原卷作「求证：$A\subset B$」（不带下划线，即真包含），本稿原写作 $\subseteq$，
%     2026-09-15 回原图核对时按原卷记号订正。
\item 若 $m>3$，求证：$A\subset B$.
\end{enumerate}

\ans{（1）$m\leqslant 1$；（2）略}

\ansspace*[7]
\end{enumerate}

\end{document}
